\part{Perspectives d'inscription des archives historiques dans l'écosystème patrimonial et scientifique de l'enseignement supérieur}
\label{part:ecosysteme-patrimonial}

Les relations entre collections, institutions et données prennent des formes différentes selon l'échelle à laquelle les archives historiques sont envisagées. Au sein de l'Université de Montpellier, elles prennent place parmi des collections bibliographiques, muséales et scientifiques qui documentent parfois les mêmes personnes, institutions, enseignements ou pratiques, tout en relevant de traditions descriptives et d'environnements documentaires distincts. L'histoire de l'enseignement médical se reconstitue en outre à partir de sources aujourd'hui réparties entre plusieurs institutions de conservation et inscrites dans différents réseaux de \gls{signalement}. Enfin, certains ensembles documentaires peuvent être sélectionnés et transformés pour constituer des \glspl{corpus-numerique} répondant à une problématique de recherche déterminée.

La première échelle est donc celle du patrimoine universitaire. Elle conduit à examiner les rapprochements susceptibles d'être établis entre archives, collections bibliographiques et objets scientifiques à partir des entités qu'ils documentent en commun. L'\gls{interoperabilite} sémantique, les \glspl{referentiel} et la formalisation des relations permettent d'envisager ces rapprochements sans uniformiser les modèles descriptifs propres à chaque type de collection.

Le changement d'échelle devient ensuite institutionnel. Les sources relatives à l'histoire de l'enseignement médical montpelliérain dépassent les limites de l'Université actuelle : leur \gls{signalement} dans des réseaux collectifs, l'identification partagée de certains de leurs acteurs et la constitution d'ensembles patrimoniaux interinstitutionnels permettent de restituer des continuités que la répartition matérielle des documents rend moins immédiatement perceptibles. À cette mise en relation des ressources succède enfin leur transformation en objets de recherche : transcription, encodage, \gls{extraction-entites}, \gls{fouille-textes} et publication des données construisent, à partir d'une sélection de sources, des \glspl{corpus-numerique} dont la structure dépend d'une question et d'une méthode scientifiques.

Cette partie suit ainsi un élargissement progressif, depuis les collections de l'Université jusqu'aux réseaux interinstitutionnels, puis des ressources patrimoniales aux données produites pour la recherche.

% ============================================================
% CHAPITRE 8
% ============================================================

\chapter{Les archives historiques dans le patrimoine de l'Université de Montpellier}
\chapitrecourt{Les archives historiques dans le patrimoine universitaire}

% ============================================================

\section{Des collections décrites dans des environnements distincts}

Les archives historiques ne constituent pas le seul ensemble patrimonial conservé et diffusé au sein de l'Université de Montpellier. Bibliothèques, collections muséales et patrimoine scientifique disposent de leurs propres dispositifs de description, de numérisation et de publication, développés selon des logiques et à des rythmes différents. Leur comparaison permet de replacer les solutions étudiées dans la partie précédente dans un environnement institutionnel plus large et d'examiner la manière dont ces différentes composantes du patrimoine universitaire sont aujourd'hui susceptibles d'être rapprochées.

\paragraph{Le patrimoine documentaire : une chaîne ancienne de conservation et de diffusion}

La Bibliothèque universitaire historique de médecine bénéficie d'une politique ancienne de conservation et de valorisation. Constituée pour l'essentiel au début du XIX\textsuperscript{e} siècle à partir des ouvrages sélectionnés par Gabriel Prunelle dans les dépôts littéraires issus des confiscations révolutionnaires, elle conserve environ cent mille imprimés antérieurs au XIX\textsuperscript{e} siècle et près de neuf cents manuscrits, dont les deux tiers sont médiévaux\footcite{Dulieu1981}\footcite{Lorblanchet2015}. Le musée Atger, formé d'environ mille dessins et cinq mille estampes donnés entre 1813 et 1832, prolonge cette conception humaniste de la formation médicale en associant sciences et arts\footcite{Lorblanchet2015}.

La visibilité de ces collections résulte d'opérations successives de conservation, de catalogage et de reproduction. Dès les années 1990, l'atelier de conservation et le service photographique assuraient le microfilmage de sécurité des manuscrits, la reproduction des dessins du musée Atger et la fourniture d'images aux chercheurs et aux éditeurs\footcite[42--43]{Nicq1996}. Le projet \emph{Cantor et musicus} associait déjà numérisation des manuscrits médiévaux, enrichissement du \gls{signalement} et élargissement de l'accès\footcite[39--44]{Vial2001}. Les réorganisations institutionnelles et les déplacements des fonds ont néanmoins fragmenté les sources relatives à leur constitution. Depuis les années 2000, l'évaluation du \gls{signalement}, l'inventaire matériel, les constats d'état et la poursuite des numérisations ont contribué à leur réappropriation patrimoniale\footcite{Denton2025}.

Cette histoire explique l'existence d'une chaîne d'accès structurée. Primo constitue l'interface publique de recherche des bibliothèques de l'Université ; les imprimés et fonds anciens sont également signalés dans le Sudoc, tandis que les manuscrits et archives conservés par la bibliothèque sont décrits dans Calames. Foli@ donne accès à une sélection de manuscrits médiévaux et d'imprimés anciens numérisés, auxquels renvoient, selon les documents, les notices des catalogues nationaux\footcite{UMServicesPatrimoine}\footcite{SCDIFolia}. En amont, les ateliers du \gls{scdi} assurent les opérations de conservation-restauration, de photographie et de numérisation nécessaires à cette diffusion\footcite{SCDI2022}. Le \gls{signalement}, la conservation matérielle et l'accès aux reproductions reposent ainsi sur plusieurs services et outils complémentaires.

\paragraph{Les collections muséales et scientifiques : une professionnalisation progressive}

Les collections muséales et scientifiques ont suivi une autre trajectoire. À l'ancienne Université Montpellier 2, leur gestion s'est professionnalisée avec le recrutement d'un premier conservateur en 2004 et la création d'un Service des collections, devenu Pôle patrimoine scientifique\footcite{Gomel2011}. Celui-ci a repris des ensembles jusque-là principalement suivis par les enseignants, chercheurs et personnels techniques\footcite{Gomel2011}.

L'inventaire associe description matérielle et scientifique, historique et photographies des objets\footcite[54--57]{BourgadeTheronAumasson2014}. Au Conservatoire d'anatomie, sa reprise doit notamment établir une numérotation stable malgré les lacunes des anciens registres\footcite{Wenger2025}. La présence de restes humains impose en outre des procédures particulières d'identification, de conservation et de présentation\footcite{Ducourau2016}. Numérisation, expositions, manifestations et actions pédagogiques donnent parallèlement accès à des collections dont certaines sont dépourvues d'espace permanent d'exposition\footcite[58--61]{BourgadeTheronAumasson2014}.

La protection au titre des monuments historiques a également soutenu leur connaissance et leur sauvegarde\footcite{Palouzie2011}. Les 5\,688 pièces du Conservatoire d'anatomie ont été classées avec leur mobilier, leurs registres, leurs planches et dessins anatomiques, leurs écorchés et leurs squelettes d'animaux\footcite{POPAnatomieMontpellier}. Leur description est accessible dans Palissy, base nationale du patrimoine mobilier du ministère de la Culture, aujourd'hui diffusée au sein de la Plateforme ouverte du patrimoine (POP). L'accueil des collections Delmas-Orfila-Rouvière et Spitzner provenant de l'Université Paris-Descartes a ensuite nécessité des moyens humains et des réserves supplémentaires, ainsi que la poursuite des inventaires, restaurations, prêts et expositions\footcite[p.~62--65]{Girard2014}\footcite{PalouzieDucourau2017}\footcite{DucourauPinail2021}..

Palissy n'en fournit cependant pas un catalogue exhaustif, puisque son périmètre dépend de la protection des objets. À l'échelle institutionnelle, l'Université déploie désormais Flora comme logiciel commun de gestion des collections muséales et scientifiques. Il doit centraliser notices, \glspl{referentiel} et images, puis permettre la publication d'une sélection dans un portail public\footnote{\emph{Présentation de l'offre Flora Software à l'Université de Montpellier}, audition de négociation, 18 juillet 2024 ; \emph{Réunion de lancement du projet}, Flora Software -- Université de Montpellier, 14 février 2025}.

\paragraph{Les archives universitaires : une reconnaissance institutionnelle plus récente}

La trajectoire institutionnelle des archives historiques doit être replacée dans celle des archives universitaires françaises, dont la prise en charge spécialisée s'est développée plus tardivement que celle d'autres composantes du patrimoine universitaire. En l'absence de services internes dédiés, leur sauvegarde a longtemps reposé sur les Archives nationales et départementales ainsi que, dans certains cas, sur des bibliothèques municipales ou universitaires\footcite{MadayMechine2011}. À une période de faible investissement des universités jusqu'aux années 1980 ont succédé une sensibilisation progressive durant la décennie suivante, puis la multiplication des services spécialisés à partir des années 2000 : trois des soixante-quinze universités recensées en disposaient en octobre 2006, contre vingt-deux à la fin de 2012\footcite{Mercier2014}.

Ce décalage a été caractérisé comme un \enquote{double déficit de légitimité}, au regard des autres patrimoines reconnus et par comparaison avec les universités étrangères\footcite{Hottin2008}. Il affecte également l'exploitation scientifique des fonds et contribue aux déséquilibres documentaires et historiographiques entre établissements parisiens et universités de province\footcite{Dubois2022}. Le cas montpelliérain s'inscrit dans cette chronologie générale tout en présentant la particularité, mise en évidence dans la première partie, d'une conservation ancienne des archives de la Faculté précédant largement la constitution de la Mission qui en assure aujourd'hui le traitement spécialisé.

\paragraph{La \gls{dcsph} : un cadre institutionnel transversal}

La fusion des universités Montpellier 1 et Montpellier 2 en 2015 a rapproché des compétences patrimoniales développées selon des orientations complémentaires : l'ancienne Montpellier 2 avait notamment structuré l'inventaire et le conditionnement des collections scientifiques ainsi que leurs relations avec l'enseignement et la recherche, tandis que Montpellier 1 avait développé la restauration et la valorisation de son patrimoine historique et de ses collections\footcite{UMSchemaDirecteurPatrimoine2025}. La création de la \gls{dcsph} dans ce contexte a réuni les fonctions de préservation du patrimoine historique et de médiation scientifique au sein d'une même direction.

Son organisation actuelle comprend cinq services, dont trois directement consacrés au patrimoine : le Service du patrimoine historique, chargé notamment de la conservation, de l'étude, de l'inventaire et de la valorisation des collections mobilières et scientifiques ; le Service des collections patrimoniales documentaires, créé en 2021, qui prend en charge les collections patrimoniales des bibliothèques universitaires ; et la Mission Archives historiques (voir annexe~\ref{annexe:dcSPH}, fig.~\ref{fig:dcSPH-2025})\footcite{UMSchemaDirecteurPatrimoine2025}. Cette organisation place ainsi dans un même cadre institutionnel les objets, les collections documentaires et les archives qui témoignent de l'histoire de l'Université.

Le \emph{Schéma directeur du patrimoine historique et des collections} de 2025 donne une portée plus large à ce rapprochement. Il rattache explicitement les collections universitaires à leurs fonctions historiques d'enseignement et de recherche et fait de la réactivation de ces liens l'un des enjeux de la politique patrimoniale de l'établissement. Dans le cas de la Faculté de médecine, cette approche est particulièrement nette : le bâtiment historique est présenté comme un point d'entrée dans l'histoire de l'enseignement médical à travers les archives, le patrimoine documentaire et les collections mobilières, réunis dans une même lecture historique\footcite{UMSchemaDirecteurPatrimoine2025}.

\paragraph{Le \gls{scdi} : des compétences mutualisées à l'échelle interuniversitaire}

À ce cadre interne à l'Université s'ajoute un dispositif de coopération entre établissements. Créé le 1\textsuperscript{er} janvier 2021 dans le prolongement de missions auparavant exercées par la Bibliothèque interuniversitaire puis réparties entre les services documentaires issus des recompositions universitaires montpelliéraines, le \gls{scdi} intervient pour les universités de Montpellier dans des domaines qui gagnent à demeurer mutualisés : informatique documentaire, patrimoine écrit et graphique, circulation des documents, coordination du Sudoc, ressources électroniques et coopération régionale ou nationale\footcite{LorblanchetArabesques111}\footcite{SCDI2022}.

Dans le domaine patrimonial, cette mutualisation repose notamment sur un atelier de conservation-restauration et un atelier de numérisation-photographie. Le premier assure évaluations sanitaires, constats d'état, dépoussiérage, conditionnement, conservation préventive et traitements de restauration ; le second réalise les prises de vue et la numérisation des documents patrimoniaux, le traitement des fichiers et leur préparation pour la diffusion\footcite{LorblanchetAteliers2024}. Leur articulation est particulièrement importante lorsqu'une reproduction suppose une préparation matérielle ou des conditions particulières de manipulation. Au cours du stage, le diagnostic de plusieurs parchemins destinés à la restauration a ainsi montré que les choix de numérisation et de diffusion ne peuvent être dissociés de l'état matériel des documents.

Le \gls{scdi} ne se substitue ni aux établissements propriétaires ni aux services responsables des collections sur lesquelles il intervient. Son intérêt pour les archives historiques tient précisément à cette répartition : la Mission conserve la connaissance archivistique et la responsabilité de ses fonds, tandis que certaines opérations spécialisées de conservation, de photographie, de numérisation ou de diffusion peuvent s'appuyer sur des compétences et des équipements déjà mutualisés. Quelques registres des archives historiques sont actuellement accessibles sous forme numérisée dans Foli@ ; à la date de rédaction de ce mémoire, il n'est toutefois pas prévu que les futures numérisations de la Mission soient diffusées par cette plateforme. La mutualisation des compétences et des équipements du \gls{scdi} ne préjuge donc pas de l'environnement retenu pour la mise à disposition des reproductions. Le fonctionnement du \gls{scdi} fournit ainsi, parallèlement à l'organisation patrimoniale portée par la \gls{dcsph}, un exemple concret d'une coopération qui ne suppose pas la centralisation matérielle ou documentaire des collections.

Archives, livres et objets apparaissent ainsi comme des traces complémentaires de l'histoire universitaire, susceptibles d'être mobilisées conjointement pour documenter les enseignements, les pratiques scientifiques et les acteurs qui les ont produits ou utilisés. Ce rapprochement institutionnel et intellectuel conduit à déplacer la question vers les données elles-mêmes : certaines de ces relations peuvent-elles être formalisées de manière à être retrouvées et exploitées au-delà des publications ou projets de médiation qui les reconstituent ponctuellement ?


% ============================================================

\section{La mise en relation des archives, des livres et des objets}

\subsection{L'\gls{interoperabilite} sémantique des collections}

\paragraph{D'un rapprochement éditorial à une mise en relation des données}

La possibilité de rapprocher des ressources conservées dans plusieurs ensembles apparaît déjà dans les pratiques de médiation des universités. La série consacrée à François Rabelais sur \emph{Ex Bibliotheca}, carnet du \gls{scdi} hébergé sur Hypothèses, construit ainsi une même enquête à partir de documents conservés aux archives historiques et à la bibliothèque universitaire historique de médecine\footcite{ExBibliotheca}. L'inscription de Rabelais en 1530 est documentée par le registre des matricules S~19 et par le \emph{Liber procuratorum} S~2 ; son doctorat de 1537 par le registre des actes S~5 ; son enseignement des \emph{Pronostics} d'Hippocrate par le registre des leçons S~4\footcite{LorblanchetRabelais1}\footcite{LorblanchetRabelais2}. Ces sources archivistiques sont rapprochées des œuvres et éditions conservées par la bibliothèque, parmi lesquelles l'édition de 1532 des textes d'Hippocrate et de Galien publiée par Rabelais ainsi que différentes éditions de ses œuvres littéraires\footcite{LorblanchetRabelais4}. Guillaume Rondelet peut de même être suivi d'une ressource à l'autre : condisciple et ami de Rabelais, il apparaît dans le \emph{Tiers Livre} sous la figure de Rondibilis, tandis que ses propres ouvrages sont conservés dans les collections bibliographiques\footcite{LorblanchetRabelais3}.

Le rapprochement repose ici sur le travail intellectuel qui identifie les personnes, les événements, les œuvres et les documents pertinents puis les réunit dans un même récit. Il fournit un exemple concret de relations déjà établies entre les collections sans être exprimées comme telles dans leurs données descriptives. Leur formalisation permettrait de rendre une partie de cette contextualisation exploitable indépendamment de la publication qui l'a produite.


\paragraph{Du Web de documents au Web de données}

Archives, bibliothèques, musées et collections scientifiques reposent sur des traditions descriptives différentes. Un instrument de recherche restitue la provenance et l'organisation hiérarchique d'un fonds ; un catalogue de bibliothèque décrit des ressources bibliographiques ; un inventaire muséal individualise les objets, leur matérialité, leur provenance et leur état. Dans le cas du patrimoine scientifique, ces différences documentaires se superposent à des statuts juridiques et à des usages eux-mêmes variables\footcite{Cornu2016}.

La publication de ces \glspl{metadonnees} sur le Web ne suffit pas à assurer leur \gls{interoperabilite}. Le \gls{web-semantique} vise à rendre les données et les relations qu'elles entretiennent interprétables par les machines. Les \glspl{donnees-liees}, ou \emph{Linked Data}, en constituent une mise en œuvre fondée sur l'identification, la publication et la connexion de données structurées sur le Web\footcite{BizerHeathBernersLee2009}. Les ressources sont identifiées par des \glspl{uri} accessibles par \gls{http} et peuvent être décrites au moyen de \gls{rdf} (\emph{Resource Description Framework}), modèle standardisé de représentation des données sous la forme de relations entre ressources. Ces relations permettent notamment de relier une ressource à d'autres entités identifiées dans le même jeu de données ou dans des jeux de données extérieurs. Le passage d'un Web de documents à un Web de données consiste ainsi à rendre identifiables et manipulables non seulement des pages ou des notices, mais les entités auxquelles elles se rapportent et les relations qui les unissent.

Cette évolution est particulièrement intéressante pour des collections décrites dans des systèmes différents. Une personne, une institution, un lieu, une œuvre ou un événement peuvent être représentés comme des entités distinctes des notices dans lesquelles ils apparaissent et servir de points de jonction entre plusieurs jeux de données. Les \glspl{referentiel} jouent alors un rôle important : ils stabilisent l'identification de certaines entités et permettent de rapprocher des données produites indépendamment. Dans les institutions patrimoniales, cette logique peut contribuer à mettre en relation archives, bibliothèques et musées sans imposer la fusion de leurs systèmes descriptifs\footcite{Marcondes2012LinkedData}.


\paragraph{Du modèle conceptuel au graphe \gls{rdf}}

Une telle mise en relation suppose d'abord de déterminer ce que l'on souhaite représenter. La modélisation conceptuelle distingue les classes d'entités pertinentes pour un domaine, leurs caractéristiques et les relations qu'elles peuvent entretenir. Une personne, une activité, une œuvre ou un document constituent ainsi des catégories différentes, dont les occurrences concrètes sont des instances. Pour représenter informatiquement un ensemble dans lequel l'analyse porte principalement sur les relations entre ces entités, une organisation sous la forme d'un graphe est particulièrement adaptée : les entités y constituent des nœuds reliés entre eux par les relations qu'elles entretiennent.

Le \gls{rdf}, ou \emph{Resource Description Framework}, fournit un modèle standardisé pour représenter ce graphe. L'information y est décomposée en \glspl{triplet-rdf} associant un sujet, un prédicat et un objet\footcite{W3CRDF11Concepts}. Le sujet est une ressource identifiée par une \gls{uri} ; le prédicat, également identifié par une \gls{uri}, exprime une propriété ou la nature d'une relation ; l'objet peut être une autre ressource identifiée par une \gls{uri} ou un littéral, par exemple un nom, une date ou une valeur numérique. Les \glspl{triplet-rdf} qui partagent les mêmes ressources forment progressivement un graphe.

Cette structure permet de représenter aussi bien les caractéristiques d'une entité que ses relations avec d'autres entités. Une propriété reliant une ressource à un littéral peut par exemple exprimer une date, tandis qu'une propriété reliant deux ressources établit une relation entre elles. Le rattachement d'une instance à une classe est lui-même exprimé en \gls{rdf} par la propriété \verb|rdf:type|. Les classes et les propriétés peuvent être définies dans des vocabulaires ou des \glspl{ontologie} afin que leur signification soit documentée et réutilisable au-delà du jeu de données qui les emploie.

Le dossier documentaire constitué autour de Rabelais permet d'en proposer une représentation volontairement simplifiée. Le doctorat du 22 mai 1537 et l'enseignement des \emph{Pronostics} d'Hippocrate sont documentés respectivement par le registre des actes S~5 et le registre des leçons S~4\footcite{LorblanchetRabelais2}. L'édition lyonnaise de 1532 conservée sous la cote J~346 rassemble notamment les \emph{Pronostics} parmi les textes d'Hippocrate et de Galien publiés par Rabelais à partir des textes qu'il avait commentés à Montpellier\footcite{LorblanchetRabelais4}. Les \emph{Pronostics} peuvent dès lors être représentés comme une entité commune à l'activité d'enseignement attestée par les archives et à la ressource bibliographique conservée par la bibliothèque.

Avant de recourir à un vocabulaire formel, ces relations peuvent être décomposées sous la forme sujet--prédicat--objet afin d'expliciter la structure logique des assertions :

\begin{table}[htbp]
	\centering
	\small
	\begin{tabular}{p{0.26\textwidth} p{0.27\textwidth} p{0.35\textwidth}}
		\toprule
		\textbf{Sujet} & \textbf{Prédicat} & \textbf{Objet} \\
		\midrule
		
		François Rabelais
		& obtient
		& Doctorat du 22 mai 1537 \\
		
		Doctorat du 22 mai 1537
		& est documenté par
		& S~5, fol.~33 \\
		
		François Rabelais
		& dispense
		& Enseignement des \emph{Pronostics} \\
		
		Enseignement des \emph{Pronostics}
		& est documenté par
		& S~4, fol.~36 \\
		
		Enseignement des \emph{Pronostics}
		& porte sur
		& \emph{Pronostics} d'Hippocrate \\
		
		François Rabelais
		& publie
		& \emph{Hippocratis ac Galeni libri aliquot} \\
		
		\emph{Hippocratis ac Galeni libri aliquot}
		& contient
		& \emph{Pronostics} d'Hippocrate \\
		
		\bottomrule
	\end{tabular}
	
	\caption[Décomposition de relations documentaires autour de François Rabelais]
	{Décomposition simplifiée de plusieurs relations documentaires autour de François Rabelais sous la forme sujet--prédicat--objet. Les prédicats sont employés ici en langage naturel pour expliciter la structure des assertions avant leur traduction dans un vocabulaire formel.}
	\label{tab:rabelais-triplets-rdf}
\end{table}

Le tableau~\ref{tab:rabelais-triplets-rdf} explicite la structure logique des assertions sans constituer encore un jeu de données \gls{rdf} publiable. Dans une représentation effective en \gls{rdf}, les ressources devraient être identifiées par des \glspl{uri} et les propriétés empruntées à des vocabulaires existants ou définies dans un modèle documenté.


\paragraph{Du graphe à sa sérialisation en Turtle}

Le \gls{rdf} ne doit pas être confondu avec une syntaxe particulière d'encodage : il définit un modèle de données en graphe, qui peut être écrit selon plusieurs syntaxes de sérialisation. Turtle en constitue une représentation textuelle conçue pour être relativement compacte et lisible par un humain\footcite{W3CTurtle2014}. Les préfixes y permettent d'abréger les \glspl{uri} répétées ; le point termine un ensemble d'assertions et le point-virgule permet d'associer plusieurs prédicats à un même sujet sans répéter son \gls{uri}.

Afin de montrer concrètement la forme que pourraient prendre les assertions précédentes, l'exemple suivant utilise l'espace de noms fictif \verb|https://example.org/um/|. Les \glspl{uri} ainsi formées n'existent pas dans les systèmes de l'Université de Montpellier et servent uniquement à illustrer la syntaxe Turtle ; elles ne préjugent donc d'aucune politique institutionnelle d'identification des ressources.

\begin{verbatim}
	@prefix ex:  <https://example.org/um/> .
	@prefix xsd: <http://www.w3.org/2001/XMLSchema#> .
	
	ex:Rabelais
	ex:obtient ex:Doctorat1537 ;
	ex:dispense ex:EnseignementPronostics ;
	ex:publie ex:J346 .
	
	ex:Doctorat1537
	ex:date "1537-05-22"^^xsd:date ;
	ex:estDocumentePar ex:S5_fol33 .
	
	ex:EnseignementPronostics
	ex:estDocumentePar ex:S4_fol36 ;
	ex:porteSur ex:PronosticsHippocrate .
	
	ex:J346
	ex:contient ex:PronosticsHippocrate .
\end{verbatim}

Le préfixe \verb|ex:| constitue ici une simple abréviation : \verb|ex:Rabelais| correspondrait par exemple à l'\gls{uri} fictive \verb|https://example.org/um/Rabelais|. La déclaration \verb|@prefix| évite ainsi de répéter l'adresse complète devant chaque ressource. L'exemple distingue également deux formes de propriétés. \verb|ex:date| relie une ressource à un littéral typé : l'indication \verb|xsd:date| précise que la chaîne \verb|1537-05-22| doit être interprétée comme une date. À l'inverse, \verb|ex:porteSur|, \verb|ex:contient| ou \verb|ex:estDocumentePar| relient deux ressources identifiées. L'identifiant fictif \verb|ex:PronosticsHippocrate| fonctionne ainsi comme un nœud partagé par plusieurs branches du graphe.

Cet exemple ne constitue pas un jeu de données publiable en l'état. Les prédicats \verb|ex:obtient|, \verb|ex:dispense|, \verb|ex:publie|, \verb|ex:porteSur| ou \verb|ex:estDocumentePar| ont été créés pour rendre immédiatement lisibles les relations observées dans le corpus et ne sont définis par aucun vocabulaire partagé. Une mise en œuvre effective demanderait de déterminer quelles entités disposent déjà d'identifiants dans des \glspl{referentiel} externes, d'attribuer des \glspl{uri} aux ressources locales devant être publiées, de procéder lorsque cela est pertinent à leur \gls{reconciliation} avec des \glspl{referentiel} existants et de remplacer ces prédicats démonstratifs par des propriétés dont la sémantique est formellement définie.


\paragraph{De l'expérimentation à une infrastructure institutionnelle}

À l'échelle de l'Université, une telle mise en œuvre supposerait d'abord la définition d'une politique institutionnelle d'identification. Un espace d'\glspl{uri} placé sous un domaine contrôlé durablement par l'établissement pourrait identifier les ressources qui ne disposent pas déjà d'un identifiant externe approprié : objets scientifiques, activités, documents ou autres entités locales selon le périmètre retenu. Cette politique devrait déterminer les catégories d'entités auxquelles l'Université attribue ses propres identifiants, leurs règles de construction, l'autorité responsable de leur attribution et les conditions de leur maintien dans le temps. Une \gls{uri} publiée ne devrait notamment pas être réattribuée à une autre ressource si celle qu'elle identifie disparaît d'une application. La stabilité de ces identifiants permettrait ainsi de dissocier l'identification des entités de l'évolution des logiciels dans lesquels leurs descriptions sont produites.

La constitution du graphe n'impliquerait pas de remplacer les systèmes documentaires existants ni de saisir une seconde fois l'ensemble de leurs données. Les descriptions continueraient à être produites dans les environnements adaptés aux différentes collections. Des traitements d'extraction sélectionneraient dans ces sources les informations nécessaires au graphe, par exemple au moyen d'un export ou d'une \gls{api} lorsque le système en propose une. Des traitements de transformation convertiraient ensuite leur structure et leurs valeurs vers les classes et propriétés retenues pour le modèle commun. À ces opérations pourrait s'ajouter l'\gls{alignement}, qui établit des correspondances entre les éléments de modèles ou de vocabulaires différents, ainsi que la \gls{reconciliation}, qui cherche à déterminer si plusieurs occurrences désignent une même entité. Un nom de personne provenant de plusieurs systèmes pourrait ainsi être rapproché d'une même autorité IdRef, tandis que les informations propres à chaque description resteraient administrées dans leur environnement d'origine.

Les \glspl{triplet-rdf} ainsi produits pourraient être conservés dans un \emph{triple store}, logiciel ou système de gestion de données spécialisé dans le stockage et l'interrogation de graphes \gls{rdf}. Il existe à cette fin des solutions libres susceptibles d'être déployées sur une infrastructure administrée par l'établissement, comme Apache Jena, dont le serveur Fuseki permet notamment l'exposition et l'interrogation de données \gls{rdf}, ainsi que des solutions commerciales ou des services hébergés. Le recours à un \emph{triple store} ne supposerait donc pas que l'Université développe elle-même un système de stockage. Celui-ci ne constituerait pas davantage un nouvel outil de catalogage pour les responsables des collections : il formerait une couche transversale alimentée à partir des différents systèmes documentaires et consacrée à la mise en relation de certaines de leurs données.

L'exploitation de cette couche pourrait reposer sur plusieurs modes d'accès. Une \gls{api} permet à une application d'interroger automatiquement un service et d'en récupérer les données selon des modalités définies, sans passer par une interface de consultation humaine. Un \emph{endpoint} \gls{sparql} offrirait un accès spécifiquement destiné à l'interrogation du graphe : \gls{sparql} est un langage de requête permettant de rechercher des motifs de relations dans des données \gls{rdf}, tandis que l'\emph{endpoint} constitue l'adresse du service auquel ces requêtes sont envoyées et qui en retourne les résultats. Une requête pourrait par exemple rechercher les personnes reliées à la fois à des ressources archivistiques et bibliographiques, ou les objets scientifiques associés à une activité ou à une personne également documentée par les archives. Un tel \emph{endpoint} ne devrait pas nécessairement être ouvert sans restriction au public : il pourrait être réservé à certains usages tandis qu'une \gls{api} ou des interfaces de consultation exposeraient des services plus encadrés.

Les \glspl{uri} publiées devraient enfin être résolvables sur le Web. L'accès à l'\gls{uri} d'une entité pourrait fournir une représentation adaptée à l'usage qui en est fait : une page lisible dans un navigateur pour la consultation humaine ou une représentation structurée, par exemple en Turtle ou en \gls{jsonld}, pour une application. L'\gls{uri} constituerait ainsi un point d'accès stable à l'entité indépendamment de la forme sous laquelle les informations qui la concernent sont restituées. Stockage du graphe dans un \emph{triple store}, interrogation par \gls{sparql}, exposition de services par une \gls{api} et résolution des \glspl{uri} correspondent dès lors à des fonctions distinctes, susceptibles d'être assurées par plusieurs composants techniques.

Appliquée à Montpellier, cette architecture ferait du graphe une couche de mise en relation plutôt qu'un système documentaire supplémentaire destiné à absorber les autres. Les archives historiques pourraient continuer à produire leurs descriptions archivistiques dans leur environnement propre, tandis que les données bibliographiques et celles relatives aux collections scientifiques resteraient administrées dans leurs systèmes respectifs. Des traitements automatisés pourraient extraire les informations nécessaires depuis ces différents environnements, les transformer selon le modèle commun et les rapprocher, lorsque cela est pertinent, de \glspl{referentiel} partagés. Seules les entités et les relations utiles aux rapprochements retenus alimenteraient ainsi le graphe \gls{rdf}.

La mise en œuvre d'un tel dispositif ne supposerait donc pas le développement intégral d'une nouvelle infrastructure par l'Université. Elle nécessiterait en revanche de définir une politique institutionnelle pour l'attribution et la pérennité des \glspl{uri}, de déterminer le modèle et les vocabulaires partagés, d'organiser les traitements d'extraction, de transformation, d'\gls{alignement} et de \gls{reconciliation}, puis d'attribuer les responsabilités relatives à l'hébergement, à l'alimentation et à la maintenance de cette couche commune. Sa pérennité dépendrait ainsi autant de cette gouvernance que des composants techniques retenus.

\paragraph{\Glspl{referentiel}, \gls{reconciliation} et \glspl{ontologie}}

L'identification commune des entités constitue l'une des conditions de la mise en relation de jeux de données distincts. Les \glspl{referentiel} de personnes, collectivités, lieux, œuvres ou concepts peuvent servir de points d'appui à cette identification : plutôt que de rapprocher des chaînes de caractères identiques ou voisines, la \gls{reconciliation} cherche à déterminer si les données renvoient effectivement à une même entité. Un identifiant partagé peut alors jouer un rôle de pivot entre plusieurs systèmes.

Le cas de Rabelais permet d'en préciser l'intérêt à Montpellier. Son nom apparaît dans les descriptions des documents conservés aux archives historiques, dans les notices bibliographiques des ouvrages conservés par la Bibliothèque universitaire historique de médecine et dans les contenus éditoriaux qui rapprochent aujourd'hui ces ressources. La reconnaissance de la chaîne de caractères « François Rabelais » dans plusieurs bases ne suffit cependant pas à établir informatiquement qu'elles désignent une même personne. Une opération de \gls{reconciliation} consisterait à rapprocher ces différentes occurrences d'une même entité identifiée dans un \gls{referentiel} de personnes. Les descriptions resteraient administrées dans leurs environnements respectifs, mais pourraient partager ou aligner l'identifiant correspondant à Rabelais et rendre ainsi explicite leur point commun.

La même méthode pourrait être appliquée à d'autres entités attestées dans plusieurs collections de l'Université. Une collectivité ayant produit des archives et publié des documents, un professeur apparaissant à la fois comme producteur d'un fonds et auteur d'ouvrages, ou une œuvre conservée sous plusieurs formes pourraient être rapprochés à condition d'établir qu'il s'agit effectivement de la même entité. Le travail préalable consisterait donc à repérer les entités communes, à examiner les informations permettant de les identifier et à rechercher leur éventuelle présence dans des \glspl{referentiel} existants. La \gls{reconciliation} ne créerait pas la relation historique entre les ressources : elle fournirait un moyen de rendre exploitable dans les données une identification préalablement établie.

Les relations établies entre \glspl{referentiel} doivent elles-mêmes être qualifiées. Une relation d'identité entre deux ressources n'a pas la même portée qu'une correspondance entre deux concepts appartenant à des thésaurus distincts. Dans le contexte montpelliérain, l'objectif ne serait donc pas de constituer immédiatement un \gls{referentiel} unique pour l'ensemble des collections, mais de déterminer quels identifiants existants peuvent servir de points de jonction et quelles entités nécessiteraient une identification locale. L'\gls{interoperabilite} sémantique peut ainsi s'appuyer sur des données administrées dans plusieurs systèmes sans imposer leur fusion.

Les \glspl{ontologie} interviennent à un autre niveau. Elles ne servent plus seulement à déterminer que deux données renvoient à une même entité, mais à définir la nature des entités et des relations qui les unissent. Les prédicats fictifs employés dans l'exemple Turtle précédent --- \verb|ex:publie|, \verb|ex:porteSur| ou \verb|ex:estDocumentePar| --- rendent immédiatement lisibles les relations observées dans le corpus Rabelais, mais leur sémantique n'est définie par aucun modèle partagé. Une propriété formellement définie peut au contraire préciser les classes susceptibles d'apparaître comme sujet et comme objet ainsi que le sens exact de la relation. Le passage des identifiants partagés à la formalisation de ces relations conduit ainsi à examiner les \glspl{ontologie} susceptibles de fournir un modèle commun aux différentes collections ; le \gls{cidoccrm} offre dans ce domaine un cadre particulièrement adapté au patrimoine culturel.


\paragraph{\gls{cidoccrm} comme modèle conceptuel commun}

Le \gls{cidoccrm} fournit précisément un cadre conceptuel pour représenter des informations patrimoniales hétérogènes. Issu des travaux du Comité international pour la documentation de l'ICOM et normalisé par ISO~21127, il définit un ensemble de classes et de propriétés destiné à faciliter l'intégration, la médiation et l'échange d'informations entre sources documentaires\footcite{CIDOCCRM713}. Son intérêt pour des collections hétérogènes tient notamment à une modélisation orientée vers les événements : les personnes, groupes, objets, documents, lieux et autres entités persistantes peuvent être reliés par les événements et activités auxquels ils participent.

Cette organisation diffère d'une description fondée principalement sur l'accumulation d'attributs attachés à une notice. Les informations sont exprimées à travers des instances de classes et les propriétés qui les relient ; chaque propriété possède un domaine et un co-domaine qui précisent les catégories d'entités entre lesquelles elle peut être employée. Une même activité peut ainsi être reliée à un acteur, à une temporalité, aux documents qui permettent de la connaître et aux ressources qu'elle mobilise, tandis qu'un objet matériel peut être distingué du contenu symbolique qu'il porte.

Le cas de Rabelais permet d'en éprouver concrètement le principe. Dans la proposition de modélisation présentée en annexe~\ref{annexe:rabelais-cidoccrm}, François Rabelais est représenté comme une instance de \emph{E21 Person}. La collation de son doctorat et son enseignement sont représentés comme des instances de \emph{E7 Activity}, tandis que leurs temporalités relèvent de \emph{E52 Time-Span}. Les propriétés \emph{P11 had participant}, \emph{P14 carried out by} et \emph{P4 has time-span} permettent respectivement de représenter la participation de Rabelais à la collation de son doctorat, son rôle dans l'activité d'enseignement et la temporalité de ces activités.

La modélisation conduit également à distinguer le support matériel d'archives de l'information documentaire qu'il porte. Le registre des actes S~5 et le registre des leçons S~4 sont considérés, dans leur matérialité, comme des instances de \emph{E22 Human-Made Object}. Les mentions du doctorat et de l'enseignement inscrites dans ces registres sont représentées comme des instances de \emph{E31 Document}. La propriété \emph{P128 carries} relie alors le support matériel au contenu documentaire qu'il porte, tandis que \emph{P70 documents} relie ce contenu à l'activité qu'il documente. Cette distinction évite de confondre le registre comme objet conservé avec l'information documentaire inscrite sur ce support.

Une distinction comparable est opérée pour la ressource bibliographique. L'exemplaire conservé sous la cote J~346 est représenté comme une instance de \emph{E22 Human-Made Object}, tandis que les \emph{Pronostics} d'Hippocrate relèvent de \emph{E33 Linguistic Object}. La propriété \emph{P128 carries} permet de relier l'exemplaire matériel au contenu linguistique qu'il porte. L'activité d'enseignement peut en outre être reliée aux \emph{Pronostics} par \emph{P16 used specific object}, propriété qui associe une activité à la ressource spécifique qu'elle mobilise. Le rapprochement initialement exprimé en langage naturel par « l'enseignement porte sur les \emph{Pronostics} » est ainsi reformulé au moyen d'une propriété dont la sémantique est définie dans le \gls{cidoccrm}.

Le passage du tableau~\ref{tab:rabelais-triplets-rdf} et de l'exemple Turtle au schéma présenté en annexe~\ref{annexe:rabelais-cidoccrm} montre ainsi la différence entre la structure d'un graphe et sa modélisation sémantique. Le \gls{rdf} permet d'exprimer des assertions sous la forme de \glspl{triplet-rdf}, mais n'impose pas à lui seul la signification de propriétés telles que « documente », « participe à » ou « porte ». Le recours au \gls{cidoccrm} contraint au contraire leur emploi par la définition des classes, des propriétés, de leurs domaines et de leurs co-domaines. La formalisation ne consiste donc pas à transformer mécaniquement les verbes d'une analyse documentaire en prédicats, mais à déterminer si les relations observées peuvent être exprimées par les concepts d'un modèle partagé.

La proposition demeure nécessairement partielle. Le choix des classes et des propriétés dépend du niveau de granularité retenu et des informations effectivement attestées par les sources ; une modélisation plus poussée pourrait notamment distinguer plus finement l'édition de 1532, l'exemplaire J~346 et les différents textes qu'il contient. Son intérêt réside ici dans la possibilité de remplacer certaines relations exprimées en langage courant par des relations dont la signification et les conditions d'emploi sont explicitement définies dans le \gls{cidoccrm}.

Une telle modélisation peut être construite à partir de descriptions existantes. Les \glspl{metadonnees} des Archives nationales portugaises ont ainsi fait l'objet d'une mise en correspondance d'éléments issus de \gls{isadg} avec les classes et propriétés du \gls{cidoccrm}. Cette transformation permet de conserver l'information archivistique tout en la rendant exploitable dans un modèle fondé sur des entités et des relations\footcite{MeloRodriguesVaragnolo2023}. La description hiérarchique du fonds conserve sa fonction archivistique, tandis qu'une représentation complémentaire autorise d'autres parcours à travers les personnes, les lieux, les événements ou les temporalités documentés.

Le Musée virtuel de l'Université de Barcelone fournit un point de comparaison particulièrement proche du contexte montpelliérain. Les collections scientifiques, artistiques et facultaires de l'établissement, auxquelles s'articulent les ressources patrimoniales de la bibliothèque et des archives, étaient initialement documentées dans plusieurs environnements ; leur rapprochement a conduit à rechercher une représentation transversale respectant la diversité de ces collections\footcite{SalseRoviraEtAl2024UniversityHeritage}. La solution retenue repose sur un \gls{profil-application} fondé sur \gls{dcmi} Metadata Terms, enrichi notamment par des éléments du \gls{cidoccrm} et de \gls{lido}, afin d'obtenir un niveau d'\gls{interoperabilite} compatible avec les moyens et les pratiques quotidiennes des responsables de collections\footcite{SalseRoviraMateoBretos2025DublinCore}.

L'expérience barcelonaise montre également que l'\gls{interoperabilite} dépend du travail effectué sur les données en amont de leur publication. Nettoyage, normalisation et \gls{reconciliation} permettent d'identifier plus sûrement les personnes, lieux et autres entités et de les rapprocher de \glspl{referentiel} externes. OpenRefine est notamment employé dans ce processus de préparation à l'enrichissement sémantique\footcite{SalseRoviraMateoBretos2025}. La \emph{Minimum Record Recommendation} poursuit une logique comparable en proposant aux petites collections un noyau descriptif réduit mis en correspondance avec le \gls{cidoccrm}, \gls{lido} et l'\gls{edm}\footcite{StadtlerStricker2024}.

À Montpellier, le corpus Rabelais fournit ainsi un premier ensemble suffisamment circonscrit pour éprouver cette démarche entre archives et collections bibliographiques. Les personnes, actes universitaires, activités d'enseignement, œuvres et éditions identifiés dans \emph{Ex Bibliotheca} permettent de tester la définition des classes utiles, le niveau de granularité des entités, la \gls{reconciliation} avec des \glspl{referentiel} existants et la sélection de propriétés correspondant réellement aux relations documentées. Une extension aux collections scientifiques devrait suivre la même méthode en partant de liens attestés entre documents et objets, plutôt que de rapprochements exclusivement nominaux ou thématiques. Le \emph{Schéma directeur du patrimoine historique et des collections}, qui considère conjointement les différentes traces de l'histoire de l'enseignement et de la recherche conservées par l'Université, fournit à cette approche un cadre institutionnel particulièrement favorable\footcite{UMSchemaDirecteurPatrimoine2025}.

% ============================================================

\subsection{Relier les archives aux collections scientifiques : l'apport de \gls{patstec}}

\paragraph{Documenter l'objet scientifique par ses sources}

Créée en 2003 par le ministère chargé de la Recherche et coordonnée par le Musée des Arts et Métiers, la Mission nationale de sauvegarde du patrimoine scientifique et technique contemporain, dite \gls{patstec}, anime un réseau d'universités, d'organismes de recherche, de musées et d'entreprises. Sa base nationale documente les instruments de recherche, d'enseignement et d'innovation et peut leur associer vidéos, témoignages, parcours de chercheurs, cahiers de laboratoire et autres ressources éclairant leur fonctionnement et leurs usages\footcite{PATSTECMission}.

Cette méthode présente un intérêt particulier pour les archives universitaires parce qu'elle envisage l'objet scientifique à partir de l'ensemble documentaire qui permet d'en restituer l'histoire. Les archives peuvent renseigner sa conception, son acquisition, ses transformations et son utilisation dans un laboratoire ou un enseignement, tandis que l'objet conserve la matérialité d'une pratique scientifique dont les dossiers fournissent d'autres traces\footcite{Joyaux2018}. La relation entre les deux ressources repose ainsi sur l'identification de l'instrument, des acteurs et des activités auxquels les documents se rapportent.

\gls{patstec} fournit de ce point de vue une réalisation concrète de la logique examinée précédemment : la notice d'un instrument peut être enrichie par des ressources complémentaires, notamment des témoignages, cahiers de laboratoire, photographies ou documents techniques qui en renseignent le fonctionnement et le contexte scientifique. La documentation de l'objet peut ainsi restituer les activités de recherche et d'enseignement auxquelles il a participé et les acteurs qui l'ont conçu, employé ou transformé.

\paragraph{Un cas montpelliérain : Émile Jeanbrau et la transfusion sanguine}

L'Université de Montpellier porte la mission \gls{patstec} Occitanie-Est\footcite{PATSTECOccitanieEst}. Le sous-fonds Émile Jeanbrau relatif à la transfusion sanguine fournit un exemple particulièrement précis de relation entre un objet scientifique et les archives qui permettent d'en restituer l'usage. Son instrument de recherche signale un appareil conservé au Conservatoire d'anatomie, tandis que plusieurs dossiers documentent les techniques de transfusion, les instruments employés, l'usage du sang citraté, des photographies réalisées en 1917 et des schémas de l'appareil\footcite[notamment 1EJ55, 1EJ59, 1EJ87, 1EJ91--96, 1EJ105, 1EJ107, 1EJ110 et 1EJ144]{BertrandDikoff2024Jeanbrau}.

L'appareil n'est pas identifiable dans la base publique de \gls{patstec} et sa datation excède le périmètre chronologique affiché par le programme national\footcite{PATSTECObjets}. Le cas permet en revanche d'envisager une transposition locale de la méthode : la notice de l'objet pourrait renvoyer vers les dossiers qui renseignent sa conception, ses caractéristiques ou ses usages, tandis que l'instrument de recherche pourrait identifier explicitement l'objet conservé au Conservatoire d'anatomie. Émile Jeanbrau, l'appareil, les activités de transfusion, les documents techniques, les photographies et les schémas formeraient ainsi un ensemble de ressources dont les relations peuvent être établies à partir des sources.

Le rapprochement diffère du cas Rabelais par la nature des ressources mises en relation. Autour de Rabelais, les relations associent principalement des activités universitaires, des actes d'archives, des œuvres et des ressources bibliographiques ; autour de Jeanbrau, elles relient un acteur, une pratique médicale, un instrument scientifique et les documents qui en renseignent l'histoire. Le principe documentaire reste identique : une relation transversale devient exploitable lorsqu'elle porte sur des entités suffisamment précisément identifiées et repose sur des sources qui permettent d'en qualifier la nature.


% ============================================================
% CHAPITRE 9
% ============================================================

% ============================================================
% CHAPITRE 9
% ============================================================

\chapter{Les archives historiques dans les réseaux et corpus interinstitutionnels}
\chapitrecourt{Réseaux et corpus interinstitutionnels}

% ============================================================

Le rapprochement des collections de l'Université constitue un premier niveau d'intégration. L'histoire de l'enseignement médical montpelliérain déborde cependant les limites de l'établissement : les documents qui permettent de la retracer sont aujourd'hui conservés par plusieurs institutions et décrits dans plusieurs réseaux. Leur réunion matérielle n'est ni possible ni souhaitable ; leur mise en relation repose sur des dispositifs qui interviennent à des niveaux différents. Le \gls{signalement} collectif rend les fonds repérables depuis plusieurs environnements documentaires ; les \glspl{referentiel} permettent d'identifier certaines des personnes et collectivités qui les traversent ; la candidature au registre \emph{Mémoire du monde} construit l'unité patrimoniale d'un ensemble dispersé ; Biblissima+ permet enfin d'envisager, à partir de ressources sélectionnées, des \glspl{corpus-numerique} organisés selon leur transmission et leur exploitation scientifique.

Le changement d'échelle est donc à la fois institutionnel et documentaire. Il ne consiste plus seulement à diffuser les archives historiques hors de la Faculté, mais à les rendre susceptibles d'être découvertes, rapprochées et exploitées avec des ressources conservées ailleurs. Dans cette perspective, la conservation \emph{in situ} ne constitue pas un obstacle : ce sont les descriptions, les \glspl{identifiant-perenne}, les reproductions et les données produites à partir des fonds qui circulent entre les réseaux.


% ============================================================

\section{Le \gls{signalement} des fonds et l'identification de leurs acteurs}

\subsection{Calames et FranceArchives : deux réseaux de \gls{signalement}}

\paragraph{Deux constructions collectives issues de traditions documentaires différentes}

Calames et FranceArchives ne constituent pas deux interfaces équivalentes dans lesquelles une même description serait simplement reproduite. Leur histoire et leur périmètre expliquent les fonctions différentes qu'ils peuvent remplir pour les archives historiques.

Calames est issu de la constitution, au cours des années 2000, d'un environnement collectif consacré aux manuscrits et aux archives conservés dans l'enseignement supérieur. Sa création prolonge notamment la rétroconversion en \gls{ead} du \emph{Catalogue général des manuscrits} et du répertoire PALME : l'interface publique ouvre en décembre 2007, puis l'outil professionnel de catalogage partagé est déployé auprès d'un premier groupe d'établissements en 2008\footcite{AbesHistoireCalames}\footcite{NoelCalames2020}. Le réseau réunit désormais une soixantaine d'établissements de l'\gls{esr}, dans lesquels peuvent être signalés manuscrits médiévaux, archives scientifiques, fonds littéraires, iconographie ou papiers de chercheurs\footcite{AbesCatalogueCalames2026}. À la fin de 2023, sa base de production comptait plus de 1,74 million de composants \gls{ead} ; 56\,970 nouveaux composants avaient été publiés au cours de cette seule année et environ 78\,000 nouveaux liens entre descriptions et autorités  avaient été ajoutés\footcite{AbesCalamesStat2023}.

Cette volumétrie donne une portée différente au \gls{signalement} envisagé pour Montpellier. L'intégration des archives historiques ne consisterait pas seulement à disposer d'une nouvelle adresse publique pour leurs instruments de recherche : elle les placerait dans un catalogue où les unités archivistiques sont mises en relation avec des fonds de même nature produits ou conservés dans d'autres établissements de l'\gls{esr}. Une recherche portant sur un professeur, une discipline, une société savante ou une institution scientifique pourrait ainsi faire apparaître conjointement les ressources montpelliéraines et des archives universitaires ou scientifiques conservées ailleurs. La description hiérarchique continuerait de restituer la provenance et la structure propre de chaque fonds, tandis que les points d'accès et les \glspl{donnee-autorite} permettraient les rapprochements transversaux.

FranceArchives procède d'une autre logique. Ouvert en 2017, le portail national agrège les instruments de recherche de services relevant de réseaux beaucoup plus divers : Archives nationales, services ministériels, Archives départementales, municipales ou métropolitaines, établissements publics, associations et autres institutions conservant des archives. En mai 2026, 241 partenaires y participaient et 108\,504 inventaires y étaient intégrés ; le portail recevait en moyenne environ 300\,000 visiteurs uniques par mois\footcite{FranceArchivesRejoindre2026}. FranceArchives constitue également l'agrégateur national vers le Portail européen des archives, de sorte que le \gls{signalement} local peut s'inscrire dans une chaîne de découverte dépassant le seul cadre français\footcite{FranceArchivesRejoindre2026}.

\paragraph{Reconstituer un paysage documentaire montpelliérain sans réunir les fonds}

Cette différence de périmètre présente un intérêt particulier à Montpellier. L'histoire institutionnelle de l'ancienne Université de médecine a produit des archives aujourd'hui réparties entre la Faculté, les Archives départementales de l'Hérault et les Archives municipales de Montpellier. La dispersion résulte de trajectoires documentaires et institutionnelles différentes ; elle ne saurait être corrigée par une réunion artificielle des descriptions dans un nouveau portail local.

Le \gls{signalement} collectif offre une autre solution. Dans Calames, les archives historiques seraient replacées dans l'histoire documentaire de l'enseignement supérieur et de la recherche ; dans FranceArchives, leurs instruments pourraient être rapprochés de ceux des autres services montpelliérains et héraultais. Une même recherche sur l'enseignement médical pourrait ainsi conduire à des fonds relevant de plusieurs institutions sans que leur provenance ni leur organisation soient effacées. Le parcours de recherche devient distribué de la même manière que les sources elles-mêmes.

Cette logique est d'autant plus importante que FranceArchives ne prétend pas constituer une reproduction exhaustive des fonds de chacun de ses partenaires : les institutions n'y versent pas nécessairement tous leurs inventaires et le portail recommande de revenir aux services producteurs lorsque la recherche n'y aboutit pas\footcite{FranceArchivesContenu2026}. Le portail fonctionne donc comme une couche nationale de découverte plutôt que comme un substitut aux systèmes locaux. Cette organisation rejoint directement l'architecture retenue dans ce mémoire : maintenir la description dans l'environnement documentaire responsable de sa production, puis l'exposer ou la réutiliser dans plusieurs interfaces.

À Montpellier, la complémentarité de Calames et FranceArchives pourrait donc être évaluée non en fonction du nombre d'interfaces dans lesquelles apparaît un instrument, mais de la capacité de chacune à l'inscrire dans un contexte documentaire différent. Un instrument maintenu dans Calames pourrait rejoindre d'autres dispositifs par les mécanismes d'exposition et de \gls{moissonnage} étudiés précédemment ; FranceArchives fournirait parallèlement le contexte archivistique national. La même description pourrait ainsi participer à plusieurs réseaux sans devenir plusieurs descriptions autonomes qu'il faudrait maintenir séparément.

Cette perspective reste suspendue aux transformations actuelles de Calames déjà examinées dans la partie précédente. La suspension temporaire des nouveaux déploiements et la préparation de son remplacement empêchent d'en faire une solution immédiatement disponible. Elles ne remettent cependant pas en cause l'intérêt du modèle collectif : quelle que soit l'architecture du futur outil, la question demeure celle d'une production structurée à l'échelle de l'\gls{esr}, réutilisable dans plusieurs contextes de découverte.


\subsection{IdRef et l'identification des personnes au-delà des collections}
\label{sec:idref}

\paragraph{Du point d'accès à l'entité identifiée}

Le rapprochement devient plus précis lorsque les points d'accès sont reliés à des \glspl{referentiel}. Dans Calames, personnes et collectivités peuvent être associées à IdRef, administré par l'\Gls{abes}. Le dispositif ne constitue plus seulement une liste contrôlée de formes de noms : il attribue des \glspl{identifiant-perenne} à des entités qui peuvent être réutilisés dans plusieurs productions documentaires de l'\gls{esr}\footcite{MistralNicolas2017}\footcite{MistralIdRef2024}.

Cette fonction a progressivement pris une ampleur importante. Fin 2024, IdRef donnait accès à environ 6,4 millions de notices d'autorité, dont près de 4,2 millions de personnes physiques et plus de 420\,000 collectivités\footcite{AbesRapportActivite2024}. Le réseau de production associe environ 180 établissements documentaires et alimente plusieurs environnements au-delà du Sudoc : Calames, les dispositifs nationaux consacrés aux thèses, Persée et différentes plateformes scientifiques\footcite{AbesReseauAutorites2026}. IdRef participe également au \gls{web-semantique} : son \gls{referentiel} est exposé en \gls{rdf} et le \emph{triple store} data.idref.fr permet son interrogation en \gls{sparql}. En 2022, celui-ci exposait environ 231 millions de \glspl{triplet-rdf}, reliant 5,4 millions de notices d'autorité à 13 millions de références bibliographiques du Sudoc et à plus de 800\,000 références issues de plusieurs portails HAL\footcite{AbesDataIdRef}.

L'intérêt pour Montpellier tient précisément à cette capacité de circulation. L'analyse des demandes adressées à la Mission a montré que les recherches prennent très fréquemment pour point d'entrée un nom de personne ou un repère chronologique, bien plus souvent qu'une cote déjà identifiée (voir fig.~\ref{fig:point-entree-demandes}). Ce constat concerne toutefois en grande partie les archives modernes et contemporaines, notamment les recherches relatives à la scolarité, et ne doit pas être projeté artificiellement sur les seuls registres anciens. Des recherches nominatives portent néanmoins aussi sur des acteurs des fonds anciens, comme Pierre Tremolet : l'enjeu commun est alors de passer d'un nom fourni par l'usager aux différentes traces documentaires susceptibles de concerner la même personne.

Une personne peut apparaître dans les registres de la Faculté comme étudiant ou gradué, dans les archives administratives comme professeur, dans les catalogues bibliographiques comme auteur, dans une thèse comme directeur ou candidat et dans un fonds privé comme producteur. La répétition de son nom dans ces différentes descriptions ne garantit pas à elle seule qu'elles désignent la même personne. Une \gls{donnee-autorite} rassemble les formes retenues et rejetées du nom, fournit des éléments de désambiguïsation et associe l'entité à un \gls{identifiant-perenne} susceptible d'être repris dans plusieurs systèmes\footcite{MistralIdRef2022}.

Le cas de François Rabelais, utilisé au chapitre précédent pour l'\gls{interoperabilite} sémantique, illustre cette différence entre chaîne de caractères et entité. Dans les archives, les catalogues et les publications de médiation, son nom apparaît dans des contextes documentaires distincts ; le rattachement de ces occurrences à un même identifiant permettrait de signaler qu'elles renvoient au même acteur sans imposer une description commune des ressources. À plus grande échelle, la même méthode pourrait être appliquée aux professeurs, étudiants et savants pour lesquels une identification suffisamment sûre peut être établie.


\paragraph{IdRef dans l'écosystème des identifiants persistants de l'\gls{esr}}

La fonction d'IdRef doit toutefois être replacée dans un écosystème plus large d'\glspl{identifiant-perenne}. Les recommandations publiées en 2025 par le Comité pour la science ouverte soulignent le rôle de ces identifiants dans la circulation des informations entre les systèmes de l'\gls{esr}, leur fiabilisation et leur réutilisation. La stratégie proposée ne repose pas sur la généralisation d'un identifiant unique à toutes les catégories d'entités, mais sur l'articulation de dispositifs adaptés à des objets et à des usages différents, notamment entre des identifiants internationaux et les \glspl{referentiel} nationaux existants\footcite{StollLamotte2025}.

Le chantier engagé par CollEx-Persée en fournit une application institutionnelle récente. Le groupement cherche à être identifié simultanément dans plusieurs dispositifs complémentaires : le \gls{ror} (\emph{Research Organization Registry}), registre international consacré à l'identification pérenne des organisations de recherche ; l'\gls{isni} (\emph{International Standard Name Identifier}), identifiant normalisé international attribué aux personnes et aux collectivités impliquées dans la création, la production, la gestion ou la diffusion de contenus ; ainsi que les \glspl{referentiel} de la Bibliothèque nationale de France et IdRef. Leur mise en relation avec Wikidata doit faciliter la circulation et la réutilisation de ces identifiants entre plusieurs systèmes\footcite{CollExPerseeIdentifiants2026}. L'exemple concerne ici une organisation plutôt qu'une personne, mais il met en évidence un principe également applicable à l'identification des acteurs : une même entité peut disposer d'identifiants dans plusieurs \glspl{referentiel} complémentaires, dont les correspondances peuvent être explicitées plutôt que remplacées par un identifiant institutionnel unique.

Cette articulation précise la fonction qu'IdRef pourrait remplir pour les archives historiques. Pour les personnes qui y sont déjà identifiées, l'identifiant IdRef peut constituer un point de jonction avec d'autres productions documentaires de l'\gls{esr} sans se substituer aux autres identifiants dont elles pourraient disposer. Pour celles qui ne figurent dans aucun \gls{referentiel} pertinent, notamment certains étudiants connus uniquement par les archives, une identification locale demeure nécessaire. L'enjeu réside donc dans la capacité à établir et à documenter les correspondances pertinentes entre identifiants plutôt que dans la constitution d'un \gls{referentiel} universel.


\paragraph{Un pivot documentaire, non un modèle prosopographique}

L'intérêt d'IdRef ne doit cependant pas conduire à faire de la \gls{donnee-autorite} le réceptacle de toutes les informations historiques relatives à une personne. Les inscriptions, grades, lieux d'origine, fonctions, carrières et relations interpersonnelles constituent des données dépendant d'une source, d'un contexte et d'une problématique de recherche\footcite{BissonGoloubkoff2020}. Leur structuration relève d'un modèle prosopographique ou historique spécifique.

Cette distinction apparaît de manière particulièrement nette dans le projet CHExRISH de l'Université Jagellonne de Cracovie. Le \emph{Corpus Academicum Cracoviense} réunit environ 67\,000 enregistrements relatifs aux étudiants et diplômés de l'Université entre 1364 et 1780 ; les personnes y sont associées à des événements académiques, professionnels et biographiques, que l'équipe cherche à intégrer avec les ressources de la bibliothèque et des autres collections universitaires au moyen du \gls{cidoccrm}\footcite{DoValleMirandaKuttNalepa2025}. L'identité d'une personne constitue ainsi un nœud du modèle, tandis que ses diplômes, fonctions ou activités demeurent représentés comme des événements distincts.

Pour une prosopographie des étudiants montpelliérains, IdRef pourrait remplir une fonction comparable de point d'ancrage extérieur. Les personnes déjà identifiées dans le \gls{referentiel} pourraient être reliées à leur \gls{identifiant-perenne} IdRef ; les autres conserveraient un identifiant interne au \gls{corpus-numerique}. Les informations issues des registres resteraient liées à leurs sources et seraient structurées dans le modèle de l'enquête. Cette répartition évite deux écueils : considérer toutes les occurrences d'un même nom comme une seule personne et transformer une notice d'autorité en biographie exhaustive.

Le lien avec l'\gls{interoperabilite} sémantique devient alors direct. IdRef répond principalement à une question d'identification : \emph{de quelle entité s'agit-il ?} Lorsqu'une même personne dispose d'identifiants dans plusieurs \glspl{referentiel}, leur \gls{alignement} peut rendre explicites les correspondances entre les systèmes qui les emploient. Un modèle comme le \gls{cidoccrm} ou une base prosopographique permet ensuite de représenter les événements, les fonctions, les lieux ou les documents auxquels cette entité est reliée. Pour les archives historiques, l'articulation entre identification partagée, \gls{alignement} des identifiants et données propres au \gls{corpus-numerique} pourrait ainsi constituer l'un des principaux moyens de passer du \gls{signalement} documentaire à l'exploitation scientifique.


% ============================================================

\section{La candidature \emph{Mémoire du monde} comme corpus patrimonial interinstitutionnel}

\paragraph{La construction d'un ensemble documentaire réparti entre plusieurs institutions}

Créé par l'\gls{unesco} en 1992, le programme \emph{Mémoire du monde} vise à favoriser la préservation, l'accès et la connaissance d'un patrimoine documentaire dont la signification peut dépasser l'institution qui le conserve\footcite{UNESCOMemoireMonde}\footcite{Grunberg2016}\footcite{MaurelMemoireMonde2020}. L'inscription au Registre international n'entraîne ni transfert de propriété ni déplacement des documents : elle reconnaît un ensemble documentaire dont les différentes composantes demeurent conservées par leurs détenteurs\footcite{UNESCORegistreMemoireMonde}\footcite{WilsonMemoryWorld2019}. Cette possibilité est explicitement prévue par la procédure de candidature. Pour le cycle 2026--2027, ouvert du 1\textsuperscript{er} juillet au 30 novembre 2025, les propositions devaient être transmises avec l'accord préalable des propriétaires ou dépositaires et porter sur un patrimoine documentaire fini et précisément délimité ; lorsque celui-ci était réparti entre plusieurs lieux ou détenteurs, chacune de ses composantes et chacun de ses responsables devaient être identifiés\footcite{UNESCOAppel2026_2027}.

La candidature montpelliéraine des \enquote{collections documentaires de l'enseignement médical à Montpellier du XII\textsuperscript{e} au XIX\textsuperscript{e} siècle}, déposée dans ce cycle, associe ainsi l'Université de Montpellier, la Ville de Montpellier, Montpellier Méditerranée Métropole et le Département de l'Hérault autour d'un ensemble défini par l'histoire de l'enseignement médical plutôt que par son lieu actuel de conservation\footcite{ParcoursCandidatureUNESCO2025}. L'unité revendiquée n'est ni celle d'un fonds au sens archivistique ni celle d'une collection matériellement constituée : elle repose sur les relations historiques entre des documents qui témoignent de la continuité des institutions, des enseignements et des pratiques médicales.

Le parcours publié à l'occasion de la candidature permet d'observer concrètement cette distribution :

\begin{landscape}
	\begin{table}[p]
		\centering
		\small
		\renewcommand{\arraystretch}{1.25}
		
		\begin{tabular}{
				@{}
				p{0.30\linewidth}
				p{0.55\linewidth}
				p{0.10\linewidth}
				@{}
			}
			\toprule
			\textbf{Institution de conservation}
			& \textbf{Ressource}
			& \textbf{Cote} \\
			\midrule
			
			Archives municipales de Montpellier
			& Transcription de l'édit de Guilhem VIII de 1181 dans le \emph{Cartulaire des Guilhem}
			& AA~1 \\
			
			Archives municipales de Montpellier
			& Bulle \emph{Quia sapientia}
			& Louvet~619 \\
			
			Archives municipales de Montpellier
			& Diplôme de bachelier de 1645
			& 15~S~5 \\
			
			\addlinespace[0.3em]
			
			Archives départementales de l'Hérault
			& Copies des statuts de Conrad
			& G~1124--G~1136 \\
			
			Archives départementales de l'Hérault
			& \emph{Privilegia Universitatis medicae Monspeliensis}
			& G~1277 \\
			
			Archives départementales de l'Hérault
			& Plan de l'École de chirurgie Saint-Côme
			& C~524-7 \\
			
			Archives départementales de l'Hérault
			& Affiche annonçant l'inauguration de l'École de santé
			& L~2498 \\
			
			\addlinespace[0.3em]
			
			Archives historiques de la Faculté de médecine
			& \emph{Liber lectionum et clavium}
			& S~4 \\
			
			Archives historiques de la Faculté de médecine
			& \emph{Liber matriculae}
			& S~19 \\
			
			Archives historiques de la Faculté de médecine
			& Registre \emph{Privilèges et statuts}
			& S~1 \\
			
			Archives historiques de la Faculté de médecine
			& Registre de délibérations
			& 1~MED~45 \\
			
			\addlinespace[0.3em]
			
			Médiathèque Émile-Zola
			& Recueil constitué par Laurent Joubert
			& Ms.~242 \\
			
			\bottomrule
		\end{tabular}
		
		\caption[Répartition des ressources du parcours \emph{Mémoire du monde}]
		{Répartition entre plusieurs institutions des ressources documentaires présentées dans le parcours publié à l'occasion de la candidature montpelliéraine au programme \emph{Mémoire du monde}.}
		\label{tab:memoire-monde-repartition}
	\end{table}
\end{landscape}

Les informations du tableau~\ref{tab:memoire-monde-repartition} sont issues du parcours publié pour accompagner la candidature\footcite{ParcoursCandidatureUNESCO2025}. La Bibliothèque universitaire historique de médecine apporte parallèlement à l'ensemble environ 12\,000 ouvrages, parmi lesquels 513 manuscrits, près de 300 incunables et post-incunables et plus de 3\,000 thèses anciennes\footcite{ParcoursCandidatureUNESCO2025}. Les institutions ne conservent donc pas simplement différentes portions d'un même ensemble : leurs ressources documentent des aspects complémentaires de son histoire, depuis la fondation et les privilèges de l'Université jusqu'à son organisation, ses enseignements, ses étudiants, ses grades, la pratique médicale et ses transformations institutionnelles.


\paragraph{De la dispersion documentaire à une coopération interinstitutionnelle}

La préparation de la candidature impose de rendre explicites des relations que la répartition actuelle des documents ne permet pas toujours de percevoir immédiatement. La délimitation de l'ensemble suppose d'identifier les pièces qui en relèvent, leur institution de conservation, leur contribution à sa valeur patrimoniale ainsi que les conditions de leur conservation et de leur accessibilité. Cette opération revêt une importance particulière à Montpellier, où un même texte ou un même épisode institutionnel peut être documenté par une pièce originale conservée dans un service, par une copie inscrite dans un registre conservé dans un autre et par une édition imprimée accessible ailleurs. Le dossier permet de contextualiser conjointement ces ressources sans effacer leurs statuts documentaires ni leurs contextes de conservation.

Le précédent des archives de l'ancienne Université de Louvain, inscrites au Registre international en 2013, montre qu'un patrimoine universitaire peut également être reconnu lorsque ses composantes sont conservées par plusieurs détenteurs\footcite{UNESCOLouvain}\footcite{ArchivesEtatLouvain2014}. Dans le cas montpelliérain, la candidature fournit ainsi un cadre institutionnel commun à des ressources jusqu'alors administrées et décrites séparément. Elle rejoint à ce titre l'un des enjeux de ce mémoire : rendre perceptibles les relations entre des documents dispersés sans subordonner leur rapprochement à leur réunion matérielle ou à leur gestion dans un même système.

L'intérêt attendu d'une inscription dépasse la reconnaissance symbolique. L'\gls{unesco} associe au Registre des objectifs de sensibilisation, de préservation et d'accès et souligne que cette reconnaissance peut accroître la visibilité publique des collections et soutenir les démarches de financement consacrées à leur préservation\footcite{UNESCOAppel2026_2027}. Ces perspectives rejoignent plusieurs opérations identifiées dans le \emph{Schéma directeur du patrimoine historique et des collections} de l'Université : amélioration du \gls{signalement}, numérisation, mise en relation des collections, médiation et coopération entre institutions\footcite{UMSchemaDirecteurPatrimoine2025}. La candidature fournit toutefois un périmètre patrimonial, un récit commun et un cadre partenarial, non l'\gls{architecture-documentaire} nécessaire à l'exploitation transversale des ressources. Celle-ci suppose des instruments permettant de les identifier, de documenter leurs correspondances et, selon les usages, d'en extraire de nouveaux ensembles de données.


\paragraph{Du périmètre patrimonial aux corpus scientifiques}

Le dossier de candidature prévoit déjà plusieurs prolongements scientifiques et documentaires : la transcription et la traduction du \emph{Cartulaire}, une étude prosopographique des étudiants et de leurs circulations, l'identification de nouveaux sujets de recherche, la rédaction d'un guide coordonné des sources et, à terme, la réalisation d'un portail mutualisé\footnote{\emph{Collections documentaires de l'enseignement médical à Montpellier du XII\textsuperscript{e} au XIX\textsuperscript{e} siècle}, dossier de candidature au Registre international \emph{Mémoire du monde} de l'\gls{unesco}, version déposée en 2025}. Ces orientations ne déterminent cependant ni les modèles de données ni les méthodes selon lesquelles les sources seraient constituées en \glspl{corpus-numerique}. Le périmètre patrimonial rassemble des documents en raison de leur contribution à l'histoire de l'enseignement médical montpelliérain ; un \gls{corpus-numerique} résulte pour sa part d'une sélection opérée en fonction d'une question de recherche et d'une transformation des sources selon un modèle adapté à l'enquête.

L'étude prosopographique prévue par la candidature pourrait ainsi donner lieu à la constitution d'un \gls{corpus-numerique} à partir des séries documentant les étudiants et leurs parcours. Il faudrait alors déterminer les séries et les unités d'observation retenues et structurer les informations relatives aux noms, lieux d'origine, dates, inscriptions, grades et fonctions. Les \glspl{identifiant-perenne} d'IdRef évoqués précédemment pourraient être associés aux personnes déjà présentes dans le \gls{referentiel}, tandis que les autres individus conserveraient des identifiants internes au projet. Les relations entre personnes et événements seraient établies à partir des registres en maintenant, pour chaque information, une référence précise au passage qui la documente. IdRef assurerait ainsi l'identification et, lorsqu'elle est possible, la mise en relation avec d'autres ressources, tandis que la structure des données prosopographiques dépendrait des sources et des questions de l'enquête.

La transcription et la traduction du \emph{Cartulaire}, également prévues dans le dossier, pourraient être prolongées par un travail sur la transmission documentaire des actes. La mise en correspondance des textes publiés par Alexandre Germain avec les originaux subsistants, leurs copies dans les registres, leurs cotes actuelles et leurs éventuelles reproductions permettrait d'en restituer les différents états documentaires. Une table de concordance pourrait indiquer, pour chaque acte, son état de transmission, l'institution conservant le témoin, sa cote, son édition et, lorsqu'elle existe, sa reproduction numérique. Un même acte deviendrait alors navigable à la fois comme texte juridique, pièce originale, copie transmise par un registre et unité d'une édition savante.

Ces deux prolongements illustrent la différence entre la définition d'un patrimoine commun et sa transformation en objet de recherche. La candidature établit le périmètre interinstitutionnel et formule les orientations scientifiques ; la prosopographie, les concordances, les identifiants partagés et la constitution de \glspl{corpus-numerique} correspondent aux méthodes susceptibles de rendre certaines relations entre ses composantes directement exploitables pour des enquêtes déterminées.

% ============================================================

\section{Biblissima+ et la reconstitution numérique des traditions documentaires}

\paragraph{Une infrastructure de recherche fondée sur les relations entre ressources}

Biblissima+ permet de franchir une étape supplémentaire en faisant des relations entre ressources le principe même de constitution de \glspl{corpus-numerique}. L'ÉquipEx+ réunit dix-sept équipes au sein d'un consortium associant quinze établissements d'enseignement et de recherche autour de l'étude des cultures écrites anciennes et de la transmission des textes\footcite{BiblissimaPresentation2026}. Ses appels à projets soutiennent des opérations de recherche, de documentation, de numérisation et de valorisation portant sur les manuscrits, imprimés anciens et autres objets portant du texte\footcite{BiblissimaPresentation2026}. L'infrastructure ne se réduit donc pas à son portail : elle fournit un cadre scientifique, des modèles, des outils et des compétences permettant de constituer et de relier des corpus.

Les données fédérées concernent notamment manuscrits, imprimés anciens, œuvres, personnes, institutions, provenances et collections historiques\footcite{BiblissimaSources}\footcite{TurcanVerkerk2020}. Ces entités permettent de suivre la transmission d'un texte, la circulation d'un volume, l'histoire d'une bibliothèque ou les activités des personnes qui leur sont associées. La logique rejoint celle examinée avec le \gls{cidoccrm} : l'intérêt réside moins dans la réunion de notices dans une interface unique que dans l'explicitation de relations entre des ressources issues de systèmes documentaires différents.

Montpellier est déjà présent dans cet environnement. Plusieurs centaines de manuscrits de la Bibliothèque universitaire historique de médecine sont signalés dans Biblissima+, à partir de données provenant notamment des ressources spécialisées fédérées par l'infrastructure\footcite{BiblissimaMontpellierMedecine}\footcite{BiblissimaRessources}. Lorsque des reproductions sont accessibles selon \gls{iiif}, elles peuvent être associées à ces données. Les travaux de reconstitution de bibliothèques dispersées, comme la bibliothèque Bouhier, montrent parallèlement comment provenances et possesseurs successifs peuvent servir de relations pour rapprocher des volumes aujourd'hui conservés dans plusieurs établissements\footcite{BiblissimaBouhier}\footcite{TurcanVerkerk2020}.

\paragraph{De l'agrégation de ressources à l'étude de leur transmission}

Pour les archives historiques, l'intérêt de Biblissima+ se situe donc dans des ensembles déterminés plutôt que dans le \gls{signalement} général du fonds. Bulles, lettres patentes, statuts et privilèges présentent notamment une histoire de transmission qui dépasse leur conservation actuelle : certains textes subsistent sous forme originale, d'autres sous forme de copies intégrées à des registres, tandis qu'une partie est connue à travers les éditions du \emph{Cartulaire de l'Université de Montpellier}.

Le projet ECRU, consacré à l'ancienne Université de Paris, offre à cet égard un point de comparaison particulièrement proche. Porté par la Bibliothèque interuniversitaire de la Sorbonne avec l'IRHT, il vise à produire une édition numérique en \gls{xml}-\gls{tei} d'éditions critiques anciennes d'actes et de lettres concernant l'Université de Paris. Le premier ensemble traité comprend les quelque 2\,600 pages du \emph{Chartularium Universitatis Parisiensis} ; l'OCR et l'encodage semi-automatique doivent permettre une exploration par lieux, personnes, collectivités et institutions de conservation\footcite{BiblissimaECRU}. L'exemple est particulièrement intéressant pour Montpellier parce qu'il montre comment un cartulaire universitaire édité peut devenir le point de départ d'un \gls{corpus-numerique} structuré et relié à d'autres ressources.

Les développements conduits autour de Telma montrent de la même manière que l'édition électronique d'actes ne consiste pas seulement à reproduire un texte imprimé : l'encodage permet d'associer analyse diplomatique, texte, indexation, informations de transmission et, lorsque cela est possible, images de la source\footcite{BertrandTelma2016}. Le passage à l'édition numérique transforme ainsi l'édition en un ensemble de données susceptible d'être interrogé et relié.

Un projet montpelliérain pourrait partir d'un ensemble limité d'actes dont plusieurs états de transmission sont identifiables. La première étape consisterait à établir, acte par acte, les relations entre le texte, ses témoins, ses copies, ses éditions et les institutions qui les conservent. Cette opération gagnerait à précéder tout choix de plateforme : elle déterminerait les entités et relations réellement documentées et permettrait d'identifier les lacunes ou les incertitudes.

Les archives historiques apporteraient à ce travail la connaissance des pièces, de leurs cotes, de leur place dans les fonds et de l'histoire des classements ; les chercheurs établiraient les textes et les relations de transmission ; les personnes, lieux et institutions pourraient être rapprochés des \glspl{referentiel} déjà mobilisés par Biblissima+. La candidature \emph{Mémoire du monde} fournit ainsi le périmètre patrimonial interinstitutionnel, tandis qu'un projet Biblissima+ pourrait sélectionner un sous-ensemble textuel et en expliciter les relations documentaires.

Cette différence d'échelle est essentielle. Le futur portail mutualisé envisagé dans la candidature n'aurait pas nécessairement vocation à devenir l'environnement scientifique de toutes les éditions, de même que Biblissima+ n'aurait pas vocation à absorber les instruments archivistiques. Chaque dispositif répond à une fonction distincte : le premier peut organiser l'accès transversal au patrimoine montpelliérain ; les seconds permettent d'étudier plus finement certains ensembles selon une problématique et un modèle de données déterminés.


% ============================================================
% CHAPITRE 10
% ============================================================

\chapter{Les archives historiques comme corpus et données de recherche}
\chapitrecourt{Corpus et données de recherche}

% ============================================================

Le \gls{signalement} collectif et les infrastructures patrimoniales rendent les sources repérables et permettent de constituer des corpus dépassant les frontières institutionnelles. Leur exploitation scientifique engage une transformation supplémentaire : les documents sélectionnés deviennent la matière de transcriptions, d'éditions, d'index, d'annotations ou de jeux de données produits selon une question, une méthode et une chaîne de traitements déterminées.

Cette transformation dépasse la numérisation. La mise à disposition d'une image facilite l'accès à la source, mais ne la rend pas pour autant directement interrogeable par les méthodes computationnelles. Les travaux consacrés aux archives et aux \glspl{donnees-liees} soulignent précisément la différence entre la disponibilité numérique d'un document et la production de données structurées, intégrables et interrogeables par machine\footcite{Hawkins2022ArchivesLinkedData}. Pour les archives historiques, le passage du document au \gls{corpus-numerique} suppose donc de documenter les opérations qui transforment successivement l'image, le texte et les informations extraites de celui-ci.


% ============================================================

\section{CollEx-Persée et la constitution de corpus de recherche}

\paragraph{Du fonds conservé au corpus construit par la recherche}

Créé en 2017 puis devenu infrastructure nationale de recherche en information scientifique et technique en 2021, le \gls{gis} CollEx-Persée réunit établissements de l'\gls{esr}, bibliothèques, services d'archives et opérateurs nationaux autour du \gls{signalement}, de la numérisation, de l'enrichissement et de l'exploitation scientifique des collections\footcite{DesosWarnier2022}\footcite{CollExMissions}. Sa logique déplace progressivement la notion de collection d'excellence vers celle de corpus construit avec les communautés scientifiques et accompagné par des services adaptés à son exploitation\footcite{CollExFeuilleRoute2025}.

Cette évolution est particulièrement importante pour les archives historiques. Un fonds est constitué par son origine et son histoire documentaire ; un \gls{corpus-numerique} de recherche est construit selon une question. Les deux périmètres peuvent se recouvrir, mais ils ne se confondent pas. Une étude des mobilités étudiantes pourrait sélectionner plusieurs registres appartenant à des séries différentes ; une recherche sur l'histoire d'un enseignement pourrait associer délibérations, programmes, cours et imprimés ; un projet consacré à un professeur pourrait réunir son fonds, ses publications et des objets scientifiques liés à son activité.

Les projets soutenus par CollEx-Persée montrent cette variété. Le \emph{Métadictionnaire médical multilingue} structure les entrées de quarante-neuf dictionnaires et encyclopédies médicales des XVIII\textsuperscript{e} au XX\textsuperscript{e} siècle ; \emph{Digital Alfieri} fédère des manuscrits et livres annotés répartis entre plusieurs institutions ; les archives Gilliéron associent documents iconographiques, archives textuelles et objets et complètent leur description en \gls{ead} par des autorités faisant l'objet d'un \gls{alignement} avec IdRef, PeriodO et GeoNames\footcite{CollExMetadictionnaire}\footcite{CollExDigitalAlfieri}\footcite{CollExGillieron}.

Ces exemples montrent également que l'infrastructure n'impose ni une typologie documentaire ni une chaîne technique unique. Le point commun réside dans la constitution d'un objet scientifique suffisamment défini pour justifier une coopération entre détenteurs des collections, chercheurs et spécialistes des traitements numériques. À Montpellier, le \emph{Cartulaire}, les registres d'étudiants, les délibérations ou certains fonds scientifiques pourraient relever d'une telle logique, mais pour des raisons différentes et avec des traitements distincts.

\paragraph{Dimensionner un projet montpelliérain}

Le premier enjeu ne serait donc pas d'obtenir le rattachement des archives historiques à une infrastructure, mais de définir un \gls{corpus-numerique} scientifiquement justifié. Plusieurs critères devraient être réunis : une question de recherche explicite, un ensemble documentaire identifiable, un état de traitement suffisant pour localiser et citer les sources, des partenaires scientifiques, une méthode de transformation des documents et une stratégie de diffusion ou de dépôt des résultats.

Les archives scientifiques rendent particulièrement visible cette chaîne de compétences. Correspondances, carnets, photographies, données, publications, documents administratifs et objets peuvent procéder d'une même activité de recherche tout en relevant de traitements documentaires différents\footcite{Charmasson2006}. L'école d'été CollEx-Persée consacrée en juillet 2026 aux fonds, corpus et collections hétérogènes de l'enseignement supérieur a précisément abordé leur traitement intellectuel, matériel et juridique, leur \gls{interoperabilite} et leur diffusion\footcite{CampusCondorcetEcole2026}.

L'Humathèque Condorcet fournit un exemple d'organisation intégrant conservation de fonds de chercheurs, traitement archivistique, numérisation, structuration des \glspl{metadonnees} et projets de transcription\footcite{HumathequeArchives}\footcite{HumathequeRecherche}. Montpellier présente une organisation différente mais dispose de compétences susceptibles d'être articulées : la Mission connaît et décrit les fonds ; le \gls{scdi} possède des compétences en conservation-restauration, photographie, numérisation et traitement documentaire ; la \gls{dsin} intervient sur les infrastructures techniques ; les chercheurs définissent les problématiques et méthodes d'exploitation.

Cette distribution peut constituer un avantage à condition que les responsabilités soient explicites. La Mission n'aurait pas à devenir une équipe d'humanités numériques, pas davantage qu'une équipe de recherche ne devrait redécrire les fonds indépendamment des instruments existants. La première garantirait l'identification et le contexte des sources ; les chercheurs définiraient le \gls{corpus-numerique} et son modèle ; les services documentaires et techniques assureraient les opérations correspondant à leurs compétences. CollEx-Persée pourrait alors servir de cadre de coopération et de financement à un projet déterminé plutôt que constituer une nouvelle couche documentaire permanente.


% ============================================================

\section{La transformation des images en textes exploitables}

\subsection{La \gls{htr} appliquée aux séries manuscrites}

\paragraph{De l'image consultable au texte interrogeable}

La numérisation produit d'abord des images. Pour les manuscrits, leur transformation en texte exploitable informatiquement peut s'appuyer sur la \gls{htr}. Une revue systématique portant sur 381 publications consacrées à l'utilisation de Transkribus entre 2015 et 2020 montre que la technologie est désormais mobilisée dans les bibliothèques, les archives et de nombreuses recherches en sciences humaines afin d'accélérer la transcription, de permettre la recherche en texte intégral et de rendre de grands ensembles documentaires accessibles à l'analyse\footcite{NockelsGoodingAmesTerras2022}. Les applications relevant directement des humanités représentent la majorité du corpus étudié, signe que la \gls{htr} a dépassé le stade de la seule expérimentation technique\footcite{NockelsGoodingAmesTerras2022}.

Cette distinction est importante pour Montpellier. La reproduction numérique d'un registre permet sa consultation à distance ; une transcription exploitable par machine rend possible une recherche portant sur son contenu et ouvre sur d'autres traitements : extraction de personnes ou de lieux, comparaison de formulations, indexation ou construction de données prosopographiques. L'image et le texte répondent ainsi à des usages complémentaires et doivent rester reliés.

\paragraph{Une reconnaissance dépendante du corpus et de l'usage}

La \gls{htr} ne fournit cependant pas une transcription de qualité constante indépendamment des sources. Une comparaison récente de PyLaia, HTR+, IDA, TrOCR et Titan sur des corpus présentant des langues, écritures, abréviations et orthographes différentes montre que les performances varient fortement selon les caractéristiques documentaires et selon la possibilité d'entraîner ou d'ajuster les modèles\footcite{RomeinRabusLeifertStroebel2025}. Les auteurs soulignent également l'importance des modèles de langue et la difficulté d'établir une comparaison totalement indépendante des paramètres propres à chaque moteur\footcite{RomeinRabusLeifertStroebel2025}.

Ces résultats invitent à renverser la question pour les archives historiques : il ne s'agirait pas de déterminer quel outil est « le meilleur », mais de sélectionner un corpus et de mesurer si un traitement donné atteint une qualité suffisante pour l'usage scientifique envisagé. Les registres présentant une certaine régularité matérielle, une continuité chronologique et un nombre limité de mains offriraient un terrain plus pertinent qu'un ensemble artificiellement constitué de pièces de typologies et d'écritures très diverses.

Une expérimentation montpelliéraine pourrait être organisée selon le protocole suivant :

\begin{table}[htbp]
	\centering
	\small
	\begin{tabular}{p{0.20\textwidth} p{0.34\textwidth} p{0.36\textwidth}}
		\toprule
		\textbf{Étape} &
		\textbf{Éléments à observer} &
		\textbf{Fonction dans l'expérimentation} \\
		\midrule
		
		Sélection de la série
		& Régularité matérielle, nombre de mains, langues, période, continuité et qualité des images
		& Déterminer si le \gls{corpus-numerique} présente une cohérence suffisante pour entraîner ou évaluer un modèle \\
		
		Constitution d'un échantillon
		& Pages représentatives des différentes mains, mises en page et difficultés du corpus
		& Éviter d'évaluer le traitement sur les seules pages les plus régulières \\
		
		Production d'une \gls{verite-terrain}
		& Transcription manuelle contrôlée de l'échantillon
		& Fournir une référence pour l'entraînement et l'évaluation \\
		
		Évaluation
		& Calcul du \gls{taux-erreur-caracteres} et examen qualitatif des erreurs
		& Mesurer la qualité du résultat et identifier les difficultés récurrentes \\
		
		Adéquation à l'usage
		& Lecture assistée, recherche plein texte, \gls{extraction-entites} ou constitution de données
		& Déterminer si la qualité obtenue est suffisante pour la fonction scientifique envisagée \\
		
		\bottomrule
	\end{tabular}
	
	\caption[Protocole possible d'évaluation d'un traitement HTR]
	{Principales étapes proposées pour l'évaluation d'un traitement de \gls{htr} appliqué à une série manuscrite des archives historiques.}
	\label{tab:protocole-htr}
\end{table}

Le tableau~\ref{tab:protocole-htr} constitue un protocole proposé, non le résultat d'une expérimentation déjà conduite. Il souligne surtout qu'une valeur de \gls{taux-erreur-caracteres} n'a de sens qu'au regard de l'usage recherché. Une transcription destinée à assister la lecture humaine peut rester partiellement fautive ; une extraction automatique de noms, de dates ou de lieux est beaucoup plus sensible aux erreurs touchant précisément ces informations.

Cette dépendance au corpus restitue une fonction centrale à la connaissance archivistique. Le mode de production du registre, les changements de mains ou de formulaire, la langue, la chronologie et la continuité matérielle déterminent en partie la faisabilité du traitement. La sélection des sources ne précède donc pas simplement l'automatisation : elle constitue l'une des conditions de sa validité méthodologique.


% ============================================================

\subsection{La transcription structurée et l'édition numérique}

\paragraph{De la transcription à une chaîne éditoriale réutilisable}

La transcription obtenue par \gls{htr} ne constitue qu'un état intermédiaire. Le projet DAHN, développé dans le cadre d'une collaboration entre Inria, Le Mans Université et l'EHESS, propose une chaîne ouverte articulant numérisation, segmentation, transcription, post-traitement, encodage et publication, avec la \gls{tei} comme format pivot\footcite{Chiffoleau2024TEIPipeline}. Le choix d'un format structuré indépendant de l'interface vise précisément à assurer la robustesse, la pérennité et la réutilisabilité des corpus produits\footcite{Chiffoleau2024TEIPipeline}.

Dans cette chaîne, les reproductions peuvent être déposées dans Nakala et exposées selon \gls{iiif} ; segmentation et transcription sont réalisées avec Kraken et eScriptorium ; le texte est corrigé puis encodé en \gls{xml}-\gls{tei} avant d'être publié au moyen de TEI Publisher\footcite{Chiffoleau2024TEIPipeline}. L'intérêt pour Montpellier réside moins dans la reprise à l'identique de cette combinaison logicielle que dans le découplage des étapes. Une image, une transcription, un fichier structuré et une interface de publication constituent des objets distincts ; chacun peut évoluer sans imposer la reconstruction de toute la chaîne.

Cette organisation prolonge directement l'\gls{architecture-documentaire} mise en évidence dans la partie précédente. Les images des archives historiques pourraient être conservées dans l'environnement institutionnel retenu et exposées selon \gls{iiif} ; le \gls{corpus-numerique} de recherche conserverait ses transcriptions et fichiers \gls{tei} dans l'environnement du projet ; l'édition ou le mini-site pourrait restituer conjointement le texte et l'image au moyen de leurs \glspl{identifiant-perenne}.

\paragraph{Deux applications différentes : actes et registres}

Cette méthode serait particulièrement adaptée à l'édition du \emph{Cartulaire} envisagée dans le cadre de \emph{Mémoire du monde}. Chaque acte pourrait constituer une unité textuelle à laquelle seraient associés transcription, traduction, date, personnes et institutions mentionnées, informations de transmission, cotes des témoins et reproductions disponibles. L'encodage permettrait ainsi de dépasser la reproduction numérique de l'édition Germain en explicitant les relations entre le texte édité et les documents qui le transmettent.

L'exemple du projet ECRU consacré au \emph{Chartularium Universitatis Parisiensis} confirme qu'un tel traitement est adapté aux corpus universitaires anciens\footcite{BiblissimaECRU}. Le parallèle doit cependant s'arrêter au principe méthodologique : le modèle d'encodage montpelliérain devrait être construit à partir de ses propres actes, de leurs états de transmission et des objectifs scientifiques définis pour le projet.

Un corpus de registres appellerait un autre schéma. Pour des matricules, les unités pertinentes pourraient être les inscriptions et les personnes ; pour des délibérations, les séances, décisions, intervenants et matières ; pour des registres d'enseignement, les cours, enseignants, dates et textes étudiés. La \gls{tei} ne constitue donc pas une grille uniforme appliquée aux archives historiques : elle fournit un langage suffisamment extensible pour représenter les phénomènes retenus par l'édition ou l'analyse.


\paragraph{L'\gls{extraction-entites} et l'identification des acteurs}

Une fois les textes disponibles sous une forme exploitable, certaines informations peuvent être repérées automatiquement. L'\gls{extraction-entites}, couramment désignée par l'acronyme \gls{ner}, permet notamment d'identifier des mentions de personnes, lieux ou organisations. Les expériences conduites sur des collections patrimoniales montrent cependant que reconnaissance et désambiguïsation doivent être distinguées : repérer une chaîne comme nom de personne ne suffit pas à déterminer l'individu auquel elle correspond\footcite{VanHoolandEtAl2015NER}.

Cette difficulté est directement applicable aux archives universitaires. Les formes latines, variantes orthographiques, changements de titulature et homonymies compliquent l'identification des étudiants et professeurs. Une chaîne de traitement pertinente pour Montpellier pourrait donc associer \gls{htr}, \gls{extraction-entites}, validation humaine et \gls{reconciliation} avec IdRef ou d'autres \glspl{referentiel}. Les personnes absentes de ces derniers conserveraient un identifiant interne au \gls{corpus-numerique}.

Les registres d'étudiants constituent un cas particulièrement intéressant. Après transcription et contrôle, l'\gls{extraction-entites} pourrait contribuer à repérer noms, lieux d'origine et dates afin de préparer une table prosopographique. Les grades et fonctions, plus dépendants de la structure du registre et du vocabulaire institutionnel, pourraient nécessiter des règles spécifiques ou une annotation plus fortement supervisée.

Le résultat ne devrait toutefois jamais être dissocié du texte qui l'a produit. Une ligne indiquant qu'un étudiant est inscrit à Montpellier à une date donnée devrait permettre de revenir au passage de la transcription, puis à la vue de la reproduction et à la cote du registre. Cette exigence relie directement l'\gls{extraction-entites} à la \gls{provenance-donnees} examinée plus loin : l'automatisation accroît l'échelle d'exploitation des sources sans supprimer la nécessité de justifier chaque information historique.


% ============================================================

\section{L'indexation automatique et la \gls{fouille-textes}}

\subsection{Vocabulaires contrôlés et indexation semi-automatique}

\paragraph{De la transcription à la recommandation de sujets}

Le passage à des \glspl{corpus-numerique} volumineux permet d'assister certaines opérations d'indexation. Une expérience menée sur 4\,550 ordonnances manuscrites de la cité-État de Berne, datées de 1528 à 1798, offre un point de comparaison particulièrement utile parce qu'elle articule plusieurs étapes examinées ici : transcription automatique, vocabulaire contrôlé et recommandation de \glspl{metadonnees}\footcite{RomeinVeldhoenRomein2026}. Les documents ont été transcrits avec Transkribus, puis plusieurs méthodes d'Annif ont été évaluées à partir d'un vocabulaire hiérarchique exprimé en \gls{skos}, modèle permettant de représenter de manière structurée des vocabulaires contrôlés et les relations entre leurs concepts, comprenant environ 1\,830 libellés\footcite{RomeinVeldhoenRomein2026}.

L'expérience ne démontre pas qu'une indexation automatique puisse remplacer l'analyse documentaire. Elle met au contraire en évidence l'interdépendance des traitements : la qualité des recommandations varie selon les méthodes employées et les erreurs de transcription peuvent altérer les mots nécessaires à la classification\footcite{RomeinVeldhoenRomein2026}. Une transcription suffisamment correcte pour être comprise par un lecteur peut donc rester insuffisante pour une opération automatisée qui dépend de formes lexicales précises.

\paragraph{Un usage possible sur les délibérations montpelliéraines}

À Montpellier, cette méthode serait surtout pertinente pour un ensemble suffisamment continu et volumineux pour que l'indexation manuelle exhaustive représente une charge importante. Les séries de délibérations constituent une piste plus convaincante qu'un traitement indifférencié de tous les fonds : elles documentent sur la longue durée des décisions relatives aux enseignements, aux personnels, aux bâtiments, aux cérémonies, aux finances ou aux relations avec d'autres institutions.

Le premier travail consisterait cependant à construire le vocabulaire à partir du corpus et des questions de recherche, et non à appliquer mécaniquement un thésaurus extérieur. Certaines notions pourraient être rapprochées de \glspl{referentiel} existants ; d'autres relèveraient du vocabulaire historique et institutionnel propre à la Faculté. Une représentation en \gls{skos} permettrait de distinguer concepts, libellés, synonymes et relations hiérarchiques tout en conservant un vocabulaire réutilisable.

L'indexation semi-automatique pourrait ensuite proposer des sujets ou signaler des passages susceptibles de relever de certains concepts. Le rôle de l'algorithme serait donc de réduire l'espace à examiner et de rendre certaines tendances repérables, tandis que la validation humaine déterminerait les \glspl{metadonnees} effectivement retenues. Une telle expérimentation aurait également un intérêt réflexif : l'écart entre recommandations et validation permettrait d'identifier les notions difficiles à reconnaître automatiquement et d'améliorer à la fois le vocabulaire et la qualité des textes.


\subsection{La \gls{fouille-textes} comme instrument d'exploration}

\paragraph{Changer d'échelle sans dissocier les résultats de leurs sources}

Une fois les textes rendus exploitables par machine, la \gls{fouille-textes} permet de dépasser l'indexation descriptive. Recherche plein texte, fréquences, cooccurrences, \gls{extraction-entites}, classification ou modélisation thématique offrent différentes manières d'explorer des ensembles qu'une lecture strictement linéaire rend difficilement maîtrisables. Ces méthodes ne répondent cependant pas aux mêmes questions et leur pertinence dépend du contenu, de la structure et de la qualité du \gls{corpus-numerique}.

Des travaux conduits sur environ 55\,000 pages d'archives concernant les politiques relatives aux réfugiés dans les années 1970 ont ainsi combiné OCR, extraction d'informations, clustering et modélisation thématique pour repérer les sujets abordés par plusieurs organisations et leur évolution\footcite{GrantSebastian2021TopicModelling}. L'intérêt de l'approche réside moins dans la production automatique d'une interprétation que dans la possibilité d'identifier des sous-ensembles ou des évolutions vers lesquels orienter ensuite la lecture historique.

Cette articulation entre analyse distante et retour au document est particulièrement importante pour les archives historiques. Les traitements computationnels produisent des catégories, probabilités, fréquences ou regroupements qui résultent du corpus choisi et des paramètres du modèle. Ils constituent donc des résultats de recherche, non un nouveau niveau objectif de description des archives.

\paragraph{Des méthodes différentes selon les séries}

À Montpellier, les délibérations pourraient être étudiées selon leur évolution lexicale ou thématique : fréquence de certaines matières, apparition de nouveaux termes institutionnels, cooccurrences entre catégories de décisions ou concentration de certains sujets à des périodes particulières. Une telle analyse pourrait contribuer à repérer des ruptures ou des séquences qui seraient ensuite examinées dans leur contexte institutionnel.

Les registres d'étudiants posent une autre question et se prêtent moins à la modélisation thématique. Leur potentiel réside davantage dans l'extraction structurée de personnes, lieux, dates et grades, suivie d'analyses prosopographiques, chronologiques ou géographiques. Les registres d'enseignement pourraient quant à eux permettre de suivre la fréquence de certains auteurs ou textes, les transformations des matières enseignées ou la succession des enseignants.

Le choix de la méthode devrait donc intervenir en dernier. La question historique détermine le \gls{corpus-numerique} ; la matérialité et l'organisation de celui-ci déterminent les possibilités de \gls{htr} et de structuration ; la qualité de ces transformations détermine enfin ce qu'il devient raisonnable de demander à la \gls{fouille-textes}. Cette succession évite que la disponibilité d'un outil devienne à elle seule le principe de sélection des sources.

Elle maintient également la continuité avec l'analyse archivistique menée dans la première partie du mémoire. Une fréquence, un cluster ou un thème ne possède de signification historique qu'à condition de savoir quels documents composent le \gls{corpus-numerique}, pourquoi ils ont été produits, quelles lacunes affectent les séries et quelles transformations leur classement a connues. La lecture computationnelle ne dispense donc pas du contexte archivistique ; elle rend au contraire sa documentation plus nécessaire lorsque l'échelle de l'analyse augmente.


% ============================================================

\section{La publication et la réutilisation des données de recherche}

\subsection{La \gls{provenance-donnees}}

\paragraph{Documenter les transformations, et pas seulement citer la source}

Les opérations précédentes produisent des objets numériques distincts des archives : transcriptions, fichiers \gls{tei}, index, annotations, tables de concordance, bases prosopographiques, résultats d'\gls{extraction-entites} ou de \gls{fouille-textes}. Chacun résulte d'une succession de choix et de transformations qui conditionnent l'interprétation des données finales.

La \gls{provenance-donnees} ne peut donc être réduite à la présence d'une cote dans une bibliographie. Les travaux consacrés à la visualisation de la provenance des documents historiques soulignent précisément que les sources numérisées peuvent connaître des transformations portant sur leur transcription, leur contenu, leur organisation ou leur forme de représentation ; documenter ces étapes permet de rendre visibles les décisions méthodologiques et curatoriales introduites entre le document et les données exploitées\footcite{VancisinEtAl2023Provenance}.

Cette approche est particulièrement adaptée aux archives historiques. Une donnée extraite d'un registre devrait conserver la référence à l'institution, au fonds ou à la série, à la cote et à la vue ou au folio correspondant ; la transcription devrait indiquer ses règles, son auteur ou son processus de production et sa version ; la donnée normalisée devrait identifier l'opération qui l'a transformée ; le jeu de données publié devrait enfin documenter son propre modèle et son historique.

Le cas du doctorat de Rabelais utilisé précédemment permet d'en représenter les niveaux sans supposer que toutes ces productions existent déjà :

\begin{table}[htbp]
	\centering
	\small
	\begin{tabular}{p{0.21\textwidth} p{0.32\textwidth} p{0.37\textwidth}}
		\toprule
		\textbf{Niveau} &
		\textbf{Ressource existante ou production envisagée} &
		\textbf{Information de \gls{provenance-donnees} à conserver} \\
		\midrule
		
		Source existante
		& Registre des actes, S~5, fol.~33
		& Institution de conservation, fonds ou série, cote et localisation dans le registre \\
		
		Transcription à produire
		& Transcription du passage relatif au doctorat
		& Référence au document source, règles de transcription, auteur ou processus de transcription et version \\
		
		Donnée structurée à produire
		& Événement correspondant au doctorat de François Rabelais, daté du 22 mai 1537
		& Référence au passage source, règle d'extraction ou de normalisation et identifiants attribués ou réutilisés pour les entités \\
		
		Jeu de données éventuel
		& Enregistrement intégré à un \gls{corpus-numerique} prosopographique ou institutionnel
		& Version du jeu de données, licence, \gls{identifiant-perenne}, modèle de données et documentation des transformations \\
		
		\bottomrule
	\end{tabular}
	
	\caption[Exemple de chaîne de provenance d'une donnée historique]
	{Exemple proposé pour documenter les transformations conduisant d'une source archivistique à une donnée structurée.}
	\label{tab:provenance-rabelais}
\end{table}

Le tableau~\ref{tab:provenance-rabelais} distingue volontairement ce qui existe de ce qui serait produit dans le cadre d'un projet. La source archivistique et sa localisation sont attestées ; transcription, donnée structurée et jeu de données correspondent à des transformations éventuelles dont la \gls{provenance-donnees} devrait être documentée si elles étaient réalisées.

\paragraph{Préserver l'incertitude et la possibilité de correction}

La traçabilité concerne également les décisions interprétatives prises au cours de la constitution d'un \gls{corpus-numerique}. Une transcription incertaine, l'identification probable d'une personne, la normalisation d'un toponyme ou le rapprochement de plusieurs variantes sous une même entité ne devraient pas disparaître derrière la donnée normalisée. Cette exigence est particulièrement importante dans les corpus prosopographiques, dont les informations sont fréquemment incomplètes, imprécises ou contradictoires. Leur modélisation peut dès lors intégrer explicitement différents degrés et formes d'incertitude portant sur les personnes, les lieux, les événements, les dates ou les sources\footcite{AkokaComynWattiauLamasseDuMouza2019}.

Une personne pourrait ainsi être identifiée provisoirement à partir d'un nom, d'un lieu ou d'autres informations fournies par les registres, puis être rapprochée ultérieurement d'une \gls{donnee-autorite} IdRef grâce à de nouveaux éléments. L'identification initiale, les sources sur lesquelles elle repose et les modifications successives devraient rester documentées afin que le jeu de données puisse évoluer sans effacer les conditions dans lesquelles les rapprochements ont été établis. Le modèle appliqué à \emph{Studium Parisiense} fournit à cet égard un exemple particulièrement pertinent : l'incertitude peut y être attachée aux informations elles-mêmes et exprimée notamment à partir de termes présents dans les sources, tels que des indications de probabilité ou d'approximation\footcite{AkokaComynWattiauLamasseDuMouza2019}.

La même exigence de traçabilité s'applique aux données produites par \gls{htr}. La \gls{verite-terrain} dépend de conventions de transcription qui peuvent varier d'un projet à l'autre et influer sur les modèles entraînés à partir de ces données. Leur réutilisation suppose donc de documenter ces conventions, de distinguer les transcriptions manuelles, automatiques et corrigées, et, pour les transcriptions produites automatiquement, d'indiquer le modèle employé ainsi que le \gls{taux-erreur-caracteres} lorsqu'il est disponible\footcite{RomeinEtAl2024DataProvenance}. Les interventions ultérieures, notamment les corrections, l'\gls{extraction-entites}, la géolocalisation ou la mise en relation avec des \glspl{donnee-autorite}, devraient également permettre d'identifier les opérations réalisées et leurs contributeurs\footcite{RomeinEtAl2024DataProvenance}.

Ces informations relèvent de la \gls{provenance-donnees} parce qu'elles permettent de reconstituer les étapes par lesquelles une source numérisée devient une donnée susceptible d'être analysée. Une erreur de transcription peut entraîner l'absence d'une entité lors de l'\gls{extraction-entites}, puis affecter une statistique ou la représentation d'un réseau prosopographique ; conserver la trace des transformations successives permet de revenir au niveau auquel l'information a été produite ou modifiée plutôt que de considérer le résultat final comme indépendant de son processus de constitution.

La documentation du \gls{corpus-numerique} devrait ainsi réunir les critères de sélection des sources, les règles de transcription et de normalisation, le modèle de données, les versions successives, les traitements automatiques, leurs évaluations et les interventions humaines. Les données produites à partir des archives historiques pourraient alors être discutées, corrigées et réutilisées comme des résultats scientifiques dont les conditions de production et la relation avec les sources demeurent explicites.


\subsection{Les \glspl{entrepot-donnees} dans l'écosystème de science ouverte}

\paragraph{Distinguer le dépôt scientifique du réservoir patrimonial}

Ces exigences rejoignent les principes \gls{fair}, selon lesquels les données doivent être trouvables, accessibles selon des conditions explicites, interopérables et réutilisables\footcite{WilkinsonEtAl2016}. Leur application aux archives historiques concerne avant tout les productions issues des fonds : transcription, fichier \gls{tei}, table de concordance, modèle, annotations ou base prosopographique. Le document d'archives reste conservé et décrit dans son environnement patrimonial ; le jeu de données peut disposer de son propre dépôt, de son \gls{identifiant-perenne}, de ses \glspl{metadonnees} et de sa documentation.

Cette distinction est importante au regard du chapitre consacré à Nakala dans la partie précédente. Nakala peut servir de réservoir de reproductions dans certains projets, notamment lorsque celles-ci constituent elles-mêmes des données de recherche, mais son emploi n'impose pas que toutes les numérisations patrimoniales de la Mission y soient dupliquées. Son intérêt devient plus net lorsqu'un projet produit un \gls{corpus-numerique} autonome appelé à être cité et réutilisé.

Le projet DAHN fournit une illustration de cette articulation : Nakala est employé pour les ressources numérisées et leur exposition selon \gls{iiif}, tandis que le traitement conduit ensuite vers segmentation, transcription, \gls{tei} et publication\footcite{Chiffoleau2024TEIPipeline}. Les différents objets produits par la recherche restent ainsi identifiables sans être confondus avec l'interface éditoriale.

Recherche Data Gouv propose pour sa part un \gls{entrepot-donnees} national pluridisciplinaire fondé sur Dataverse ainsi qu'un catalogue destiné à rendre également repérables des données déposées dans d'autres \glspl{entrepot-donnees} français\footcite{RechercheDataGouvEntrepot}\footcite{RechercheDataGouvCatalogue}. Le choix entre Nakala, Recherche Data Gouv ou un autre \gls{entrepot-donnees} dépendrait donc moins de l'origine archivistique des documents que de la discipline, du projet, des formats produits et des infrastructures auxquelles les chercheurs sont rattachés.

ISIDORE intervient à un niveau différent : il collecte et enrichit les \glspl{metadonnees} exposées par des plateformes partenaires et contribue à la découverte des ressources en sciences humaines et sociales\footcite{HumaNumIsidore}. Il ne constitue donc pas un \gls{entrepot-donnees} supplémentaire, mais une couche de repérage susceptible d'accroître la visibilité de données déjà publiées ailleurs.

\paragraph{Une chaîne distribuée applicable aux projets montpelliérains}

Cette répartition peut être représentée à partir d'un projet hypothétique issu des archives historiques :

\begin{table}[htbp]
	\centering
	\small
	\begin{tabular}{p{0.24\textwidth} p{0.31\textwidth} p{0.35\textwidth}}
		\toprule
		\textbf{Niveau} &
		\textbf{Ressource ou environnement} &
		\textbf{Fonction principale} \\
		\midrule
		
		Source archivistique
		& Instrument de recherche des archives historiques
		& Identifier le document, son contexte de production, sa cote et sa place dans le fonds \\
		
		Reproduction numérique
		& Réservoir institutionnel ou \gls{entrepot-donnees} retenu pour le projet
		& Donner accès à l'image et, lorsque cela est pertinent, l'exposer selon \gls{iiif} \\
		
		\Gls{corpus-numerique}
		& Transcriptions, fichiers \gls{tei}, annotations, tables ou base de données
		& Transformer une sélection documentaire selon une problématique et un modèle scientifique déterminés \\
		
		Dépôt scientifique
		& Nakala, Recherche Data Gouv ou autre \gls{entrepot-donnees} adapté
		& Conserver, documenter, identifier et rendre citables les productions de la recherche \\
		
		Découverte
		& ISIDORE ou autres infrastructures d'indexation
		& Rendre les \glspl{metadonnees} du corpus repérables dans des environnements de recherche plus larges \\
		
		\bottomrule
	\end{tabular}
	
	\caption[Répartition fonctionnelle des ressources d'un corpus numérique]
	{Répartition possible des fonctions entre la source archivistique, sa reproduction, les données produites par la recherche, leur dépôt et leur découverte.}
	\label{tab:chaine-donnees-recherche}
\end{table}

Le tableau~\ref{tab:chaine-donnees-recherche} reprend à l'échelle de la recherche le principe architectural établi dans la partie précédente pour la diffusion : chaque environnement conserve la fonction pour laquelle il est le mieux adapté. L'instrument de recherche identifie et contextualise la source ; le réservoir d'images expose sa représentation numérique ; le \gls{corpus-numerique} organise les objets construits par l'enquête ; l'\gls{entrepot-donnees} assure leur conservation et leur citation ; les infrastructures de découverte accroissent leur visibilité.

Les \glspl{identifiant-perenne} et la \gls{provenance-donnees} assurent la continuité entre ces niveaux. Une personne extraite d'un registre peut être reliée à son \gls{identifiant-perenne} IdRef ; la donnée qui décrit son inscription renvoie au passage transcrit ; celui-ci renvoie à la vue de l'image ; l'image reste rattachée à l'unité archivistique et à sa cote. La circulation ne dépend donc pas d'une base réunissant toutes les représentations, mais de relations suffisamment stables pour permettre le passage de l'une à l'autre.


% ============================================================
% CONCLUSION DU CHAPITRE
% ============================================================

\section*{Conclusion du chapitre}

Les possibilités examinées dans ce chapitre dessinent plusieurs types de projets qui ne doivent pas être confondus. Les registres d'étudiants peuvent alimenter des enquêtes prosopographiques sur les origines, les mobilités et les grades ; les délibérations se prêtent davantage à l'étude des décisions et des transformations institutionnelles ; les fonds de professeurs permettent de rapprocher correspondances, cours, publications, photographies, documentation scientifique et parfois objets. La profondeur chronologique des Archives historiques autorise ainsi des \glspl{corpus-numerique} très différents, dont la méthode doit être déterminée par la question de recherche et par les caractéristiques des séries utilisées.

La chaîne numérique n'est pas davantage uniforme. Une série homogène peut justifier une expérimentation de \gls{htr} ; certains textes peuvent être structurés en \gls{tei} ; des corpus prosopographiques peuvent recourir à l'\gls{extraction-entites} et à la \gls{reconciliation} ; des ensembles volumineux peuvent être explorés par indexation semi-automatique ou \gls{fouille-textes}. La pertinence de chacune de ces opérations dépend de la qualité de celles qui la précèdent. Une donnée structurée n'est exploitable qu'à condition que son extraction, sa normalisation et sa relation avec la source soient suffisamment documentées.

Ces fonds peuvent également servir de supports à des enseignements d'histoire, d'archivistique, de patrimoine et d'humanités numériques. Un \gls{corpus-numerique} limité peut donner lieu à une transcription, à un encodage en \gls{xml}-\gls{tei}, à la création ou à l'enrichissement de \glspl{donnee-autorite}, à la structuration d'un jeu de données ou à la préparation de ressources \gls{iiif}\footcite{SAAACRLPrimarySources2018}. L'intérêt pédagogique réside dans la possibilité de suivre toute la chaîne, depuis l'analyse matérielle et archivistique du document jusqu'à la constitution de données publiables, en rendant visibles les choix et les transformations intermédiaires.

La Mission conserverait dans ces projets une fonction directement liée à sa connaissance des fonds : identifier les séries pertinentes, expliquer leur organisation et leur histoire, signaler les lacunes ou les reclassements et garantir le lien entre les productions numériques et les documents conservés. Les chercheurs et responsables pédagogiques définiraient les problématiques, les modèles et les méthodes d'analyse ; les services documentaires et les infrastructures de recherche apporteraient, selon les projets, les moyens de numérisation, de calcul, de dépôt et de diffusion.

Les Archives historiques pourraient ainsi servir de point de départ à des projets dont le périmètre et les traitements varieraient selon les questions posées. Une série peut devenir un terrain d'expérimentation méthodologique, plusieurs fonds être rapprochés autour de leurs acteurs et des collections réparties entre institutions former un \gls{corpus-numerique} commun. Dans tous les cas, l'exploitation scientifique dépend de la continuité documentaire permettant de revenir des données produites aux sources dont elles sont issues.


% ============================================================
% CONCLUSION DE LA PARTIE
% ============================================================

La place des Archives historiques dans l'écosystème patrimonial et scientifique de l'enseignement supérieur se construit à plusieurs échelles. Au sein de l'Université, les archives peuvent être rapprochées des collections bibliographiques, muséales et scientifiques à partir des personnes, institutions, activités, textes ou objets qu'elles documentent en commun. Les \glspl{referentiel} et les \glspl{identifiant-perenne} facilitent l'identification de ces entités, tandis que des modèles comme le \gls{cidoccrm} permettent de formaliser certaines de leurs relations. Ces rapprochements peuvent ainsi être établis sans imposer aux différentes collections un modèle descriptif ou un environnement documentaire unique.

À l'échelle interinstitutionnelle, les réseaux de \gls{signalement} et les dispositifs d'identification partagée permettent de restituer des continuités que la dispersion matérielle des sources rend moins immédiatement perceptibles. Calames et FranceArchives inscrivent les descriptions dans des réseaux documentaires plus larges, tandis qu'IdRef peut relier certains acteurs à leurs différentes traces. La candidature \emph{Mémoire du monde} montre, selon une autre logique, qu'un ensemble patrimonial commun peut être défini à partir de ressources conservées par plusieurs détenteurs sans abolir leurs cotes, leurs classements ni leurs institutions de conservation.

La constitution de \glspl{corpus-numerique} répond à une logique différente. Les perspectives offertes par Biblissima+ et CollEx-Persée montrent que le périmètre d'un objet de recherche dépend des sources sélectionnées, de la problématique et des traitements envisagés plutôt que des seules limites d'un fonds ou d'un ensemble patrimonial. La \gls{htr}, la \gls{tei}, l'\gls{extraction-entites}, la \gls{reconciliation} ou la \gls{fouille-textes} peuvent alors produire de nouvelles représentations et de nouvelles données à partir des documents, dont la conservation et la réutilisation peuvent s'appuyer sur les \glspl{entrepot-donnees} de la science ouverte.

Ces changements d'échelle et de représentation ont en commun la nécessité de préserver les relations entre les différentes formes prises par une même ressource. La \gls{provenance-donnees} permet de maintenir le lien entre une donnée et le passage dont elle est issue, entre une transcription et sa reproduction, puis entre celle-ci, l'unité archivistique décrite et son fonds. À l'échelle des institutions, la même exigence permet à une ressource de participer à un ensemble patrimonial ou à un \gls{corpus-numerique} commun tout en conservant l'identification de son contexte documentaire et de son lieu de conservation.

L'inscription des Archives historiques dans l'enseignement supérieur et la recherche ne suppose donc pas la concentration des ressources dans un même environnement. Le fonds matériel demeure sous la responsabilité de l'institution qui le conserve, tandis que ses descriptions, ses reproductions, les entités qu'il documente et les données produites à partir de certaines séries peuvent être mises en relation avec d'autres ressources. L'enjeu réside dans l'organisation de ces relations de manière à permettre leur circulation et leur réutilisation sans perdre la provenance, le contexte ni la responsabilité documentaire.